\documentclass[13pt,a4paper,twoside]{extarticle}

\usepackage[top=3cm, bottom=3.5cm, left=3.5cm, right=2cm]{geometry}
\usepackage{fontspec}
\usepackage{polyglossia}
\setmainlanguage{vietnamese}
\setmainfont{Times New Roman}

\usepackage{setspace}
\setstretch{1.5}

\usepackage{enumitem}
\usepackage{longtable}
\usepackage{booktabs}
\usepackage{caption}
\usepackage{float}
\usepackage{fancyhdr}
\usepackage{titlesec}
\usepackage{tocloft}
\usepackage{hyperref}
\usepackage{graphicx}
\usepackage{tabularx}
\usepackage{cite}
\usepackage{listings}
\usepackage{xcolor}
\usepackage{multirow}
\usepackage{array}

% --- Evidence box ---
\usepackage{tcolorbox}
\tcbuselibrary{skins,breakable}
\newtcolorbox{evidencebox}[1][]{
  colback=blue!3!white,
  colframe=blue!50!black,
  fonttitle=\bfseries,
  title={\makebox[0pt]{\raisebox{-1pt}{\fontsize{14}{16}\selectfont\char"2611}}\quad BẰNG CHỨNG CẦN THU THẬP},
  breakable,
  before skip=10pt,
  after skip=10pt,
  #1
}
\newcommand{\evidence}[1]{%
  \begin{evidencebox}
  #1
  \end{evidencebox}
}

\hypersetup{
  colorlinks=true,
  linkcolor=black,
  urlcolor=blue,
  citecolor=black
}

% --- Chapter/section formatting ---
\titleformat{\section}{\fontsize{14}{16}\bfseries}{Chương \thesection.\ }{0em}{}
\titleformat{\subsection}{\fontsize{13}{15}\bfseries}{\thesubsection.\ }{0em}{}
\titleformat{\subsubsection}{\fontsize{13}{15}\bfseries}{\thesubsubsection.\ }{0em}{}
\titlespacing{\section}{0pt}{18pt}{12pt}
\titlespacing{\subsection}{0pt}{12pt}{6pt}
\titlespacing{\subsubsection}{0pt}{12pt}{6pt}

% Make \section act as chapter (numbered from 1)
\renewcommand{\thesection}{\arabic{section}}
\renewcommand{\thesubsection}{\thesection.\arabic{subsection}}
\renewcommand{\thesubsubsection}{\thesubsection.\arabic{subsubsection}}

% --- Page style ---
\fancypagestyle{plain}{}
\pagestyle{fancy}
\fancyhf{}
\fancyhead[LE,RO]{\fontsize{10}{12}\selectfont\rightmark}
\fancyhead[RE,LO]{\fontsize{10}{12}\selectfont\leftmark}
\fancyfoot[C]{\thepage}
\renewcommand{\headrulewidth}{0.4pt}
\setlength{\headheight}{15pt}

% --- Figure/table caption ---
\renewcommand{\thefigure}{\thesection.\arabic{figure}}
\renewcommand{\thetable}{\thesection.\arabic{table}}

% --- Code listing style ---
\lstset{
  basicstyle=\ttfamily\fontsize{11}{13}\selectfont,
  breaklines=true,
  frame=single,
  backgroundcolor=\color{gray!5},
  columns=flexible,
  showstringspaces=false
}

% --- TOC depth ---
\setcounter{tocdepth}{2}

% --- Paragraph ---
\setlength{\parskip}{0pt}
\setlength{\parindent}{1.25cm}

\begin{document}

% ════════════════════════════════════════════════════
%  BÌA CHÍNH
% ════════════════════════════════════════════════════
\pagestyle{empty}
\newgeometry{top=3cm, bottom=3.5cm, left=3.5cm, right=2cm}
\begin{center}

{\fontsize{14}{18}\bfseries ĐẠI HỌC QUỐC GIA TP. HỒ CHÍ MINH}\\[4pt]
{\fontsize{16}{20}\bfseries TRƯỜNG ĐẠI HỌC CÔNG NGHỆ THÔNG TIN}\\[4pt]
{\fontsize{16}{20}\bfseries KHOA MẠNG MÁY TÍNH VÀ TRUYỀN THÔNG}\\[40pt]

\rule{\textwidth}{1.5pt}\\[6pt]
\rule{\textwidth}{0.3pt}\\[30pt]

{\fontsize{14}{18}\bfseries HỒ NHẬT MINH}\\[40pt]

{\fontsize{16}{20}\bfseries ĐỒ ÁN CHUYÊN NGÀNH}\\[20pt]

{\fontsize{22}{28}\bfseries Thiết kế và Triển khai Quy trình}\\[6pt]
{\fontsize{22}{28}\bfseries DevSecOps Kết hợp GitOps}\\[6pt]
{\fontsize{22}{28}\bfseries cho Hệ thống Microservices trên Nền tảng AWS}\\[12pt]

{\fontsize{18}{22}\bfseries Design and Implementation of a DevSecOps Pipeline}\\[4pt]
{\fontsize{18}{22}\bfseries Combined with GitOps for a Microservices System on AWS}\\[40pt]

{\fontsize{14}{18}\bfseries KỸ SƯ NGÀNH MẠNG MÁY TÍNH VÀ TRUYỀN THÔNG}\\[40pt]

\rule{\textwidth}{0.3pt}\\[6pt]
\rule{\textwidth}{1.5pt}\\[30pt]

{\fontsize{13}{16}\bfseries TP. HỒ CHÍ MINH, 2026}

\end{center}
\restoregeometry

\newpage

% ════════════════════════════════════════════════════
%  BÌA PHỤ
% ════════════════════════════════════════════════════
\pagestyle{empty}
\newgeometry{top=3cm, bottom=3.5cm, left=3.5cm, right=2cm}
\begin{center}

{\fontsize{14}{18}\bfseries ĐẠI HỌC QUỐC GIA TP. HỒ CHÍ MINH}\\[4pt]
{\fontsize{16}{20}\bfseries TRƯỜNG ĐẠI HỌC CÔNG NGHỆ THÔNG TIN}\\[4pt]
{\fontsize{16}{20}\bfseries KHOA MẠNG MÁY TÍNH VÀ TRUYỀN THÔNG}\\[40pt]

\rule{\textwidth}{1.5pt}\\[6pt]
\rule{\textwidth}{0.3pt}\\[30pt]

{\fontsize{14}{18}\bfseries HỒ NHẬT MINH -- 23520924}\\[40pt]

{\fontsize{16}{20}\bfseries ĐỒ ÁN CHUYÊN NGÀNH}\\[20pt]

{\fontsize{22}{28}\bfseries Thiết kế và Triển khai Quy trình}\\[6pt]
{\fontsize{22}{28}\bfseries DevSecOps Kết hợp GitOps}\\[6pt]
{\fontsize{22}{28}\bfseries cho Hệ thống Microservices trên Nền tảng AWS}\\[12pt]

{\fontsize{18}{22}\bfseries Design and Implementation of a DevSecOps}\\[4pt]
{\fontsize{18}{22}\bfseries Pipeline Combined with GitOps}\\[4pt]
{\fontsize{18}{22}\bfseries for a Microservices System on AWS}\\[40pt]

{\fontsize{14}{18}\bfseries KỸ SƯ NGÀNH MẠNG MÁY TÍNH VÀ TRUYỀN THÔNG}\\[30pt]

\begin{tabular}{rl}
\textbf{GIẢNG VIÊN HƯỚNG DẪN} & \\[6pt]
\multicolumn{2}{c}{\fontsize{14}{18}\bfseries LÊ ANH TUẤN}\\[40pt]
\end{tabular}

{\fontsize{13}{16}\bfseries TP. HỒ CHÍ MINH, 2026}

\end{center}
\restoregeometry

\newpage

% ════════════════════════════════════════════════════
%  THÔNG TIN HỘI ĐỒNG (trang trống)
% ════════════════════════════════════════════════════
\thispagestyle{empty}
\begin{center}
{\fontsize{14}{18}\bfseries THÔNG TIN HỘI ĐỒNG CHẤM ĐỒ ÁN CHUYÊN NGÀNH}
\end{center}
\vfill

\begin{tabular}{@{}p{5cm} p{10cm}@{}}
\textbf{Hội đồng:} & \dotfill \\[8pt]
\textbf{Chủ tịch:} & \dotfill \\[8pt]
\textbf{Thư ký:}   & \dotfill \\[8pt]
\textbf{Ủy viên:}  & \dotfill \\[8pt]
\textbf{Ngày bảo vệ:} & \dotfill \\[8pt]
\textbf{Điểm:}     & \dotfill \\[8pt]
\end{tabular}

\newpage

% ════════════════════════════════════════════════════
%  LỜI CẢM ƠN
% ════════════════════════════════════════════════════
\thispagestyle{empty}
\addcontentsline{toc}{section}{Lời cảm ơn}
\begin{center}
{\fontsize{14}{18}\bfseries LỜI CẢM ƠN}
\end{center}
\vspace{12pt}

Trước hết, em xin gửi lời cảm ơn chân thành đến thầy Lê Anh Tuấn -- người đã trực tiếp hướng dẫn em trong suốt quá trình thực hiện đồ án này. Khi em bắt đầu, khái niệm DevSecOps và GitOps còn khá mới mẻ với em. Những buổi trao đổi với thầy không chỉ giúp em hiểu rõ bản chất của các phương pháp luận này, mà quan trọng hơn, thầy đã gợi ý cho em hướng tiếp cận thực tế: không dừng lại ở lý thuyết mà phải triển khai được trên AWS với một hệ thống microservices cụ thể. Sự sát sao của thầy -- từ việc góp ý về cấu trúc Terraform module, tổ chức Jenkins Shared Library, cho đến cách trình bày báo cáo -- là yếu tố then chốt để em hoàn thành đồ án với đầy đủ các thành phần CI/CD, security scanning, và monitoring như hiện tại.

Em cũng xin cảm ơn các thầy cô trong Khoa Mạng Máy Tính và Truyền Thông. Những môn học về mạng máy tính, bảo mật thông tin, điện toán đám mây và DevOps mà em được học trong chương trình đã cung cấp nền tảng vững chắc để em tự tin làm việc với các công nghệ trong đồ án -- từ Terraform, Kubernetes, Jenkins, đến các công cụ bảo mật như SonarQube, Trivy, OWASP Dependency Check.

Cuối cùng, em muốn cảm ơn gia đình và những người bạn đã đồng hành cùng em. Có những đêm debug Terraform state conflict đến 3 giờ sáng, hay những lần Jenkins pipeline fail liên tục không rõ nguyên nhân -- chính sự động viên và kiên nhẫn lắng nghe của mọi người đã giúp em vượt qua những lúc nản lòng nhất để hoàn thiện đồ án này.

TP. Hồ Chí Minh, ngày ... tháng ... năm 2026\\
Sinh viên thực hiện\\[20pt]
{\it Hồ Nhật Minh}

\newpage

% ════════════════════════════════════════════════════
%  MỤC LỤC
% ════════════════════════════════════════════════════
\pagestyle{empty}
\addcontentsline{toc}{section}{Mục lục}
\renewcommand{\contentsname}{\centerline{\fontsize{14}{18}\bfseries MỤC LỤC}}
\tableofcontents
\newpage

% ════════════════════════════════════════════════════
%  DANH MỤC HÌNH
% ════════════════════════════════════════════════════
\thispagestyle{empty}
\addcontentsline{toc}{section}{Danh mục hình}
\renewcommand{\listfigurename}{\centerline{\fontsize{14}{18}\bfseries DANH MỤC HÌNH}}
\listoffigures
\newpage

% ════════════════════════════════════════════════════
%  DANH MỤC BẢNG
% ════════════════════════════════════════════════════
\thispagestyle{empty}
\addcontentsline{toc}{section}{Danh mục bảng}
\renewcommand{\listtablename}{\centerline{\fontsize{14}{18}\bfseries DANH MỤC BẢNG}}
\listoftables
\newpage

% ════════════════════════════════════════════════════
%  DANH MỤC TỪ VIẾT TẮT
% ════════════════════════════════════════════════════
\thispagestyle{empty}
\addcontentsline{toc}{section}{Danh mục từ viết tắt}
\section*{\centerline{\fontsize{14}{18}\bfseries DANH MỤC TỪ VIẾT TẮT}}

\begin{longtable}{@{}p{3.5cm} p{12cm}@{}}
\textbf{Từ viết tắt} & \textbf{Nghĩa tiếng Anh / Giải thích} \\ \midrule
ADOT   & AWS Distro for OpenTelemetry \\
ALB    & Application Load Balancer \\
ALB    & Application Load Balancer \\
AMI    & Amazon Machine Image \\
ANSI   & American National Standards Institute \\
AWS    & Amazon Web Services \\
CD     & Continuous Delivery / Continuous Deployment \\
CI     & Continuous Integration \\
CNI    & Container Network Interface \\
CVE    & Common Vulnerabilities and Exposures \\
DAST   & Dynamic Application Security Testing \\
DDB    & Amazon DynamoDB \\
DNS    & Domain Name System \\
EC2    & Amazon Elastic Compute Cloud \\
ECR    & Amazon Elastic Container Registry \\
EKS    & Amazon Elastic Kubernetes Service \\
ENI    & Elastic Network Interface \\
GitOps & Operations by Git (Git-based operations model) \\
IAM    & Identity and Access Management \\
IaC    & Infrastructure as Code \\
IEEE   & Institute of Electrical and Electronics Engineers \\
IRSA   & IAM Roles for Service Accounts \\
KMS    & Key Management Service \\
KLTN   & Khóa luận tốt nghiệp \\
ML     & Machine Learning \\
MLOps  & Machine Learning Operations \\
NAT    & Network Address Translation \\
NLB    & Network Load Balancer \\
NLP    & Natural Language Processing \\
OIDC   & OpenID Connect \\
OTEL   & OpenTelemetry \\
OWASP  & Open Web Application Security Project \\
RBAC   & Role-Based Access Control \\
S3     & Amazon Simple Storage Service \\
SAST   & Static Application Security Testing \\
SCA    & Software Composition Analysis \\
SG     & Security Group \\
SNS    & Amazon Simple Notification Service \\
SQS    & Amazon Simple Queue Service \\
TLS    & Transport Layer Security \\
VPC-CNI & Amazon VPC Container Network Interface \\
VPC    & Virtual Private Cloud \\
\end{longtable}

\newpage

% ════════════════════════════════════════════════════
%  TÓM TẮT (ABSTRACT)
% ════════════════════════════════════════════════════
\setcounter{page}{1}
\thispagestyle{fancy}
\addcontentsline{toc}{section}{Tóm tắt đồ án}
\section*{\centerline{\fontsize{14}{18}\bfseries TÓM TẮT ĐỒ ÁN}}

Trong bối cảnh chuyển đổi số, các tổ chức đang dịch chuyển từ kiến trúc monolithic sang kiến trúc microservices để đạt được khả năng mở rộng và tốc độ phát triển cao hơn. Việc triển khai và vận hành hàng chục microservice đặt ra thách thức lớn về bảo mật và quản lý cấu hình. Các quy trình triển khai thủ công truyền thống tiềm ẩn nguy cơ misconfiguration, lỗ hổng bảo mật không được phát hiện kịp thời, và khó khăn trong rollback. Đồ án này hướng tới giải quyết các vấn đề trên thông qua việc thiết kế và triển khai một quy trình DevSecOps kết hợp GitOps cho hệ thống microservices trên nền tảng AWS.

Quy trình được xây dựng bao gồm: (1) Quản lý hạ tầng bằng Infrastructure as Code với Terraform (module hóa VPC, EKS, ECR với 3 môi trường dev/staging/prod); (2) Tích hợp bảo mật ``shift-left'' vào CI pipeline qua Jenkins với ba security gate là SonarQube (SAST), OWASP Dependency Check (SCA), và Trivy (Container Image Scan); (3) Triển khai tự động theo mô hình GitOps sử dụng Argo CD kết hợp Kustomize và Helm; (4) Chính sách bảo mật Kubernetes với Kyverno (Policy as Code) và kiểm tra bảo mật hạ tầng với Terrascan; (5) Giám sát tập trung với CloudWatch, Prometheus, Grafana và 10 alerting rules.

Toàn bộ hạ tầng được provision tự động trên AWS bao gồm VPC, EKS cluster v1.33, 3 managed node groups (2 t3.medium + 1 t3.large) và 5 ECR repositories. Ứng dụng Retail Store gồm 5 microservice (UI, Cart, Orders, Catalog, Checkout) được triển khai thành công qua pipeline CI/CD tự động sử dụng Jenkins (CI) và Argo CD (CD GitOps). Kết quả triển khai cho thấy toàn bộ 41 pods trên cluster ở trạng thái Running, bao gồm 5 microservice ứng dụng, MySQL catalog-db, các addon hệ thống (Argo CD, Jenkins, Prometheus, Grafana, Cert Manager, Kyverno, SonarQube), và Load Balancer Controller đang hoạt động ổn định. CI pipeline đã thực thi thành công các Kaniko build, vượt qua 3 security gate, và tự động cập nhật GitOps repo để kích hoạt Argo CD sync.

Đồ án chứng minh tính khả thi của mô hình DevSecOps kết hợp GitOps trong việc tự động hóa quy trình phát triển, kiểm thử bảo mật, Policy as Code và triển khai microservices, đồng thời cung cấp một tài liệu tham khảo thực tiễn cho các tổ chức muốn áp dụng phương pháp này trên AWS.

\vspace{12pt}
\noindent\textbf{Từ khóa:} DevSecOps, GitOps, Microservices, Kubernetes, AWS, Terraform, Jenkins, Argo CD, CI/CD, Infrastructure as Code, Kyverno, Policy as Code, Prometheus, Grafana, Terrascan.

\newpage

% ════════════════════════════════════════════════════
%  CHƯƠNG 1: MỞ ĐẦU
% ════════════════════════════════════════════════════
\section{Mở đầu}

\subsection{Lý do chọn đề tài}

Các tổ chức phát triển phần mềm hiện đại đang chuyển dịch từ kiến trúc monolithic sang kiến trúc microservices, trong đó ứng dụng được phân tách thành nhiều service độc lập -- mỗi service đảm nhiệm một chức năng nghiệp vụ riêng biệt, có thể được phát triển, kiểm thử, và triển khai độc lập với các service khác [13].

Kiến trúc microservices mang lại khả năng mở rộng linh hoạt và tốc độ phát triển cao, nhưng đi kèm là chi phí vận hành gia tăng đáng kể [8]. Khi số lượng service tăng lên, việc triển khai thủ công trở nên bất khả thi. Đồng thời, bề mặt tấn công (attack surface) mở rộng đáng kể khi mỗi service có thể là một điểm yếu tiềm tàng.

Hai phương pháp luận hiện đại đã xuất hiện để giải quyết các thách thức này:

\begin{description}[leftmargin=*,nosep]
  \item[DevSecOps] mở rộng DevOps truyền thống bằng cách đưa bảo mật vào mọi giai đoạn của vòng đời phát triển, thay vì chỉ kiểm tra bảo mật ở bước cuối cùng [2]. Cách tiếp cận ``shift-left security'' này giúp phát hiện lỗ hổng sớm, giảm chi phí sửa chữa và rút ngắn thời gian phản hồi (feedback loop).
  \item[GitOps] là mô hình vận hành lấy Git repository làm nguồn sự thật duy nhất (single source of truth) cho trạng thái mong muốn của toàn bộ hệ thống [4]. Một operator như Argo CD liên tục so sánh trạng thái thực tế của Kubernetes cluster với trạng thái khai báo trong Git, và tự động đồng bộ khi phát hiện khác biệt [5].
\end{description}

Mặc dù cả hai phương pháp luận này đều đã được chứng minh hiệu quả trong các tổ chức lớn như Google, Netflix, và Amazon, việc kết hợp chúng thành một quy trình thống nhất, hoàn chỉnh cho microservices trên AWS vẫn là một bài toán chưa có nhiều hướng dẫn thực tiễn, đặc biệt trong môi trường học thuật tại Việt Nam.

Đồ án này được thực hiện với mong muốn xây dựng một hệ thống minh họa toàn bộ quy trình DevSecOps kết hợp GitOps, từ khâu quản lý hạ tầng (Infrastructure as Code), đến CI/CD tự động với security gates, và triển khai GitOps trên Kubernetes. Hệ thống sẽ sử dụng toàn bộ công cụ open-source và dịch vụ AWS phổ biến, đảm bảo tính khả thi và khả năng áp dụng thực tế.

\subsection{Mục tiêu đồ án}

Mục tiêu tổng quát của đồ án là thiết kế và triển khai thành công một quy trình DevSecOps kết hợp GitOps, vận hành trên nền tảng AWS với một ứng dụng microservices thực tế. Các mục tiêu cụ thể bao gồm:

\begin{enumerate}[leftmargin=*,nosep]
  \item \textbf{Tự động hóa hạ tầng bằng Infrastructure as Code}: Toàn bộ hạ tầng AWS (VPC, EKS cluster, IAM roles, ECR repositories) được định nghĩa và quản lý bằng Terraform. State được lưu tập trung trên S3 với cơ chế lock qua DynamoDB.

  \item \textbf{Tích hợp bảo mật vào CI pipeline (Shift-Left Security)}: Xây dựng Jenkins pipeline tự động thực hiện SAST (SonarQube), SCA (OWASP Dependency Check), và Container Image Scan (Trivy) trước khi cho phép deploy. Chỉ những artifact vượt qua tất cả security gate mới được đẩy lên ECR và triển khai.

  \item \textbf{Triển khai GitOps với Argo CD}: Sử dụng mô hình 2-repo (App Repo + GitOps Repo). Argo CD tự động đồng bộ trạng thái từ GitOps Repo lên EKS cluster. Rollback được thực hiện bằng git revert.

  \item \textbf{Giám sát tập trung}: Thiết lập CloudWatch Logs cho EKS control plane, Prometheus + Grafana dashboards với 10 alerting rules cho cluster và ứng dụng.

  \item \textbf{Chứng minh tính khả thi}: Triển khai thành công toàn bộ 5 microservices của ứng dụng bán lẻ mẫu lên EKS thông qua CI/CD pipeline tự động end-to-end (Jenkins → ECR → Argo CD → EKS), thực hiện demo rollback bằng GitOps, và triển khai Policy as Code (Kyverno) cùng IaC security scanning (Terrascan).
\end{enumerate}

\subsection{Đối tượng và phạm vi đồ án}

\textbf{Đối tượng nghiên cứu:} Quy trình DevSecOps, mô hình GitOps, nền tảng Kubernetes (Amazon EKS), các công cụ CI/CD (Jenkins, Argo CD), công cụ bảo mật (SonarQube, OWASP DC, Trivy, Kyverno, Terrascan), và hạ tầng AWS (VPC, ECR, IAM, CloudWatch).

\textbf{Phạm vi:}
\begin{itemize}[leftmargin=*,nosep]
  \item Hệ thống được triển khai trên một EKS cluster với 3 managed node groups (2 t3.medium + 1 t3.large), không triển khai multi-cluster hay multi-region.
  \item Chỉ triển khai môi trường dev. Cấu trúc cho staging và prod được chuẩn bị (Kustomize overlays) nhưng không triển khai đầy đủ.
  \item CI pipeline sử dụng Jenkins với Shared Libraries; CD sử dụng Argo CD kết hợp Kustomize và Helm umbrella chart; không triển khai các công cụ thay thế như GitHub Actions, GitLab CI, hay Flux CD.
  \item Security scanning bao gồm SAST (SonarQube), SCA (OWASP DC), container scan (Trivy), và IaC scanning (Terrascan); Policy as Code với Kyverno (5 ClusterPolicies); không bao gồm DAST, penetration testing.
  \item Không triển khai Service Mesh (Istio/Linkerd).
\end{itemize}

\subsection{Phương pháp thực hiện}

Đồ án được thực hiện theo phương pháp thực nghiệm (empirical approach), bao gồm các bước:

\begin{enumerate}[leftmargin=*,nosep]
  \item \textbf{Nghiên cứu lý thuyết}: Tìm hiểu cơ sở lý luận về DevSecOps, GitOps, Kubernetes, AWS EKS thông qua tài liệu học thuật và tài liệu kỹ thuật chính thức từ AWS.
  \item \textbf{Khảo sát công cụ}: Đánh giá và lựa chọn các công cụ phù hợp cho từng giai đoạn của pipeline.
  \item \textbf{Thiết kế kiến trúc}: Vẽ sơ đồ tổng thể hệ thống, thiết kế luồng dữ liệu và cấu trúc workspace.
  \item \textbf{Triển khai hạ tầng}: Sử dụng Terraform để tự động hóa việc provision tài nguyên AWS.
  \item \textbf{Xây dựng pipeline}: Cấu hình Jenkins cho CI, Argo CD cho CD GitOps.
  \item \textbf{Kiểm thử và đánh giá}: Thực hiện kiểm thử end-to-end, đánh giá hiệu năng và tính bảo mật của pipeline.
\end{enumerate}

\newpage

% ════════════════════════════════════════════════════
%  CHƯƠNG 2: TỔNG QUAN
% ════════════════════════════════════════════════════
\section{Tổng quan}

Chương này trình bày cơ sở lý thuyết và các công nghệ nền tảng liên quan đến đồ án, bao gồm kiến trúc microservices, phương pháp DevSecOps, mô hình GitOps, nền tảng Kubernetes và các dịch vụ AWS được sử dụng.

\subsection{Kiến trúc Microservices}

\subsubsection{Từ Monolithic đến Microservices}

Trong mô hình phát triển phần mềm truyền thống, ứng dụng được xây dựng dưới dạng một khối thống nhất (monolithic) -- tất cả các chức năng được đóng gói trong cùng một codebase, build thành một artifact duy nhất và triển khai trên một hoặc một vài máy chủ [8]. Mô hình này có ưu điểm là đơn giản trong phát triển ban đầu và dễ kiểm thử tích hợp. Tuy nhiên, khi ứng dụng phát triển về quy mô, các hạn chế trở nên rõ ràng:

\begin{itemize}[leftmargin=*,nosep]
  \item \textbf{Scale không linh hoạt}: Khi cần mở rộng, toàn bộ ứng dụng phải được nhân bản, ngay cả khi chỉ một thành phần bị quá tải.
  \item \textbf{Phát triển chậm}: Nhiều team cùng làm việc trên một codebase dẫn đến xung đột thường xuyên và thời gian build kéo dài.
  \item \textbf{Công nghệ bị khóa cứng}: Toàn bộ ứng dụng phải dùng chung một ngôn ngữ lập trình và framework.
  \item \textbf{Triển khai rủi ro cao}: Một thay đổi nhỏ có thể yêu cầu triển khai lại toàn bộ ứng dụng.
\end{itemize}

Kiến trúc microservices giải quyết các hạn chế này bằng cách phân tách ứng dụng thành nhiều service nhỏ, độc lập [13]. Mỗi service:

\begin{itemize}[leftmargin=*,nosep]
  \item Đảm nhiệm một chức năng nghiệp vụ cụ thể (single responsibility).
  \item Có thể được phát triển bằng ngôn ngữ lập trình và framework riêng (polyglot).
  \item Có database riêng (database per service pattern).
  \item Có thể được triển khai, scale, và cập nhật độc lập.
  \item Giao tiếp với nhau qua API (REST, gRPC) hoặc message queue.
\end{itemize}

\subsubsection{Thách thức của Microservices}

Tuy kiến trúc microservices giải quyết được nhiều vấn đề của monolithic, nó cũng kéo theo các thách thức mới [8], [13]:

\begin{itemize}[leftmargin=*,nosep]
  \item \textbf{Độ phức tạp trong vận hành}: Quản lý hàng chục service yêu cầu công cụ orchestration mạnh mẽ, thường là Kubernetes.
  \item \textbf{Bảo mật phân tán}: Mỗi service là một điểm entry tiềm năng; số lượng điểm cần bảo vệ tăng theo số lượng service.
  \item \textbf{Triển khai phức tạp}: Cần pipeline CI/CD phức tạp hơn để quản lý việc build, test, và deploy nhiều service.
  \item \textbf{Giao tiếp liên service}: Cần cơ chế service discovery, load balancing, circuit breaking, và retry logic.
  \item \textbf{Tính nhất quán dữ liệu}: Mỗi service có database riêng dẫn đến thách thức về distributed transactions.
  \item \textbf{Giám sát và logging}: Cần hệ thống giám sát tập trung để theo dõi trạng thái của tất cả service.
\end{itemize}

Chính những thách thức này là động lực cho sự phát triển của các phương pháp luận như DevSecOps và GitOps [2], [3].

\subsection{DevSecOps}

\subsubsection{Khái niệm}

DevOps là một tập hợp các thực hành kết hợp giữa Development (phát triển) và Operations (vận hành) nhằm rút ngắn vòng đời phát triển hệ thống và cung cấp tính năng, sửa lỗi, cập nhật thường xuyên [1]. Các nguyên tắc cốt lõi của DevOps bao gồm tự động hóa, đo lường, chia sẻ kiến thức và văn hóa hợp tác giữa các team.

Hình \ref{fig:devsecops} minh họa sự tiến hóa từ DevOps truyền thống sang DevSecOps, trong đó bảo mật được tích hợp vào mọi giai đoạn.

\begin{figure}[H]
\centering
\begin{verbatim}
   DevOps truyền thống:           DevSecOps:
    ------                         -------------- 
   | PLAN  |                       | SECURE PLAN  |
   | CODE  |                       | SECURE CODE  |
   | BUILD |                       | SECURE BUILD |
   | TEST  |                       | SECURE TEST  |
   | DEPLOY|                       | SECURE DEPLOY|
   | MONITOR| (Security at end)   | SECURE MONITOR| (Security throughout)
    ------     ↑                   -------------- 
            --- 
      SECURITY CHECK
      (Last step)
\end{verbatim}
\caption{Sự tiến hóa từ DevOps sang DevSecOps}
\label{fig:devsecops}
\end{figure}

DevSecOps mở rộng DevOps bằng cách tích hợp bảo mật (security) vào mọi giai đoạn của quy trình phát triển [2]. Triết lý cốt lõi là ``shift-left security'' -- dời việc kiểm tra bảo mật sang trái trên timeline, càng sớm càng tốt. Điều này có nghĩa là:

\begin{itemize}[leftmargin=*,nosep]
  \item Thay vì kiểm tra bảo mật ở cuối pipeline (ngay trước khi deploy), bảo mật được kiểm tra liên tục trong suốt quá trình phát triển.
  \item Developer được trao quyền và trách nhiệm về bảo mật code của mình.
  \item Các security check được tự động hóa và tích hợp vào pipeline CI/CD.
\end{itemize}

\subsubsection{Các loại Security Testing trong DevSecOps}

Có bốn loại kiểm thử bảo mật chính trong quy trình DevSecOps:

\begin{table}[H]
\centering
\begin{tabular}{@{}p{2.5cm} p{5cm} p{7.5cm}@{}}
\toprule
\textbf{Phương pháp} & \textbf{Tên đầy đủ} & \textbf{Mô tả} \\
\midrule
SAST & Static Application Security Testing & Phân tích mã nguồn tĩnh -- quét source code để tìm lỗ hổng mà không cần chạy ứng dụng. Công cụ phổ biến: SonarQube, Checkmarx, Fortify. \\
\midrule
SCA  & Software Composition Analysis & Phân tích thành phần phần mềm -- kiểm tra lỗ hổng trong các thư viện và dependency bên thứ ba. Công cụ: OWASP Dependency Check, Snyk. \\
\midrule
Container Scan & Container Image Scanning & Quét Docker image để tìm CVE trong base image và các package cài đặt. Công cụ: Trivy, Clair. \\
\midrule
DAST & Dynamic Application Security Testing & Phân tích ứng dụng đang chạy -- tấn công mô phỏng để tìm lỗ hổng thực tế. Không được thực hiện trong phạm vi đồ án này. \\
\bottomrule
\end{tabular}
\caption{Bốn loại kiểm thử bảo mật trong DevSecOps}
\label{tab:security-testing}
\end{table}

\subsubsection{Các công cụ Security sử dụng trong đồ án}

\textbf{SonarQube} là một nền tảng mã nguồn mở để kiểm tra chất lượng mã nguồn liên tục [9]. Nó thực hiện SAST bằng cách phân tích mã nguồn với hàng trăm quy tắc (rules) để phát hiện bugs, code smells, và lỗ hổng bảo mật. SonarQube hỗ trợ hơn 25 ngôn ngữ lập trình bao gồm Java, TypeScript, Go -- phù hợp với sample app của đồ án. Kết quả phân tích được hiển thị trên dashboard với các chỉ số như Quality Gate status, Security Rating, Reliability Rating.

\textbf{OWASP Dependency Check} là công cụ SCA do OWASP phát triển [10]. Nó quét các file cấu hình dependency (pom.xml cho Java Maven, package.json cho Node.js, go.mod cho Go) và đối chiếu với cơ sở dữ liệu National Vulnerability Database (NVD) để xác định các dependency có CVE đã biết. Công cụ tạo báo cáo hiển thị CVSS score và mức độ nghiêm trọng của từng lỗ hổng.

\textbf{Trivy} là công cụ quét bảo mật toàn diện do Aqua Security phát triển [11], được thiết kế để quét Docker images, file system, và Git repository. Trivy có thể phát hiện lỗ hổng trong OS packages (Alpine, Debian, Amazon Linux) và language-specific packages (Java JAR, Node.js npm, Go modules). Điểm mạnh của Trivy là tốc độ quét nhanh và tích hợp dễ dàng vào CI pipeline.

\textbf{Kyverno} là công cụ Policy as Code cho Kubernetes, được phát triển bởi Nirmata và hiện là dự án CNCF Incubating [16]. Không giống như OPA/Gatekeeper sử dụng ngôn ngữ Rego, Kyverno cho phép định nghĩa chính sách dưới dạng Kubernetes-native YAML -- cùng ngôn ngữ mà developer đã quen thuộc. Kyverno hoạt động như một dynamic admission controller trong Kubernetes cluster, intercept các API request (create, update, delete) và kiểm tra chúng với các chính sách đã định nghĩa. Các chính sách có thể thực hiện validate (chặn resource không hợp lệ), mutate (tự động sửa đổi resource), và generate (tự động tạo resource bổ sung). Trong đồ án, 5 ClusterPolicy được triển khai bao gồm: yêu cầu pod labels chuẩn, cấm sử dụng tag \texttt{latest}, yêu cầu health probes, yêu cầu resource limits, và cấm chạy container với quyền root.

\textbf{Terrascan} là công cụ static code analysis cho Infrastructure as Code, do Tenable phát triển [17]. Terrascan quét các file cấu hình IaC (Terraform, Kubernetes, Helm, Kustomize, Dockerfile) để phát hiện các vi phạm chính sách bảo mật và compliance trước khi provision. Công cụ hỗ trợ hơn 500 chính sách dựa trên các tiêu chuẩn như CIS Benchmarks, NIST, PCI-DSS. Trong đồ án, Terrascan được sử dụng để quét toàn bộ Terraform modules trước khi triển khai, phát hiện 1 vi phạm mức LOW (VPC chưa bật flow logs). Song song đó, Checkov (do Bridgecrew/Palo Alto phát triển) cũng được thiết lập sẵn sàng để bổ sung thêm một lớp kiểm tra IaC security.

\subsection{GitOps}

\subsubsection{Khái niệm}

GitOps là một mô hình vận hành cho Kubernetes và các nền tảng cloud-native, được Weaveworks đề xuất lần đầu vào năm 2017 [4]. Nguyên lý cốt lõi của GitOps là sử dụng Git repository làm ``single source of truth'' (nguồn sự thật duy nhất) cho trạng thái mong muốn của toàn bộ hệ thống [3].

Hình \ref{fig:gitops} mô tả luồng hoạt động của GitOps.

\begin{figure}[H]
\centering
\begin{verbatim}
   Git Repo (Desired State)
            ↑
            |  git push
            |
    -------- -------- 
   |   CI Pipeline    | (Jenkins)
   |   Build → Test   |
   |   → Push Image   |
   |   → Update Repo  |
    ----------------- 
            |
            v
    ----------------- 
   |  GitOps Operator | (Argo CD)
   |   -------------  |
   |  | Pull (periodic)|- Compares desired vs actual state
   |  | Diff         | |
   |  | Sync         |- Applies changes to match desired state
   |   -------------  |
    -------- -------- 
            |
            v
   Kubernetes Cluster (Actual State)
\end{verbatim}
\caption{Mô hình hoạt động của GitOps}
\label{fig:gitops}
\end{figure}

GitOps có ba nguyên tắc chính:

\begin{enumerate}[leftmargin=*,nosep]
  \item \textbf{Khai báo (Declarative)}: Toàn bộ hệ thống phải được mô tả dưới dạng khai báo (declarative configuration), ví dụ như Kubernetes manifests.
  \item \textbf{Phiên bản hóa (Versioned)}: Trạng thái mong muốn được lưu trong Git, tận dụng các tính năng của Git như versioning, branching, pull request review.
  \item \textbf{Tự động đồng bộ (Automated Sync)}: Một operator (software agent) chạy trong cluster liên tục so sánh trạng thái thực tế với trạng thái khai báo trong Git, và tự động đồng bộ (sync) hoặc cảnh báo khi phát hiện sự khác biệt (drift).
\end{enumerate}

\subsubsection{Lợi ích của GitOps}

\begin{itemize}[leftmargin=*,nosep]
  \item \textbf{Audit trail hoàn chỉnh}: Mọi thay đổi đối với hệ thống đều có commit tương ứng trong Git. Có thể dễ dàng truy vết ai đã thay đổi gì, khi nào.
  \item \textbf{Rollback đơn giản}: Rollback = git revert. Không cần quy trình phức tạp, không cần quyền đặc biệt.
  \item \textbf{Bảo mật}: Chỉ cần cấp quyền truy cập Kubernetes cho Argo CD, developer không cần quyền kubectl trực tiếp lên cluster.
  \item \textbf{Tính nhất quán}: Toàn bộ cluster được đồng bộ về đúng trạng thái khai báo trong Git, loại bỏ configuration drift.
  \item \textbf{Tự phục hồi}: Nếu có ai đó thay đổi thủ công (kubectl edit), Argo CD sẽ tự động hoàn nguyên về trạng thái trong Git.
\end{itemize}

\subsubsection{Argo CD}

Argo CD là một công cụ GitOps khai báo (declarative GitOps tool) cho Kubernetes, hiện là dự án CNCF Incubating [5]. Argo CD hoạt động theo mô hình pull-based:

\begin{enumerate}[leftmargin=*,nosep]
  \item Argo CD được cài đặt trong Kubernetes cluster.
  \item Người dùng định nghĩa Application resource -- một Custom Resource Definition (CRD) của Argo CD -- chỉ định Git repo, path đến manifests, và cluster đích.
  \item Argo CD định kỳ (mặc định mỗi 3 phút) hoặc theo webhook pull Git repo, so sánh với trạng thái thực tế của cluster.
  \item Nếu phát hiện sự khác biệt, Argo CD hiển thị trạng thái OutOfSync trên dashboard.
  \item Có thể cấu hình auto-sync để Argo CD tự động áp dụng thay đổi.
\end{enumerate}

Argo CD hỗ trợ nhiều định dạng manifest: plain YAML, Kustomize, Helm, Ksonnet, và Jsonnet. Đồ án này sử dụng Kustomize vì tính linh hoạt và là công cụ được tích hợp sẵn trong kubectl.

\subsubsection{Hai mô hình triển khai GitOps}

Có hai chiến lược triển khai GitOps phổ biến:

\begin{enumerate}[leftmargin=*,nosep]
  \item \textbf{Single Repo}: Chỉ có một repository chứa cả mã nguồn ứng dụng và manifest triển khai. Phù hợp với dự án nhỏ, ít service.
  \item \textbf{Dual Repo} (được sử dụng trong đồ án): Hai repository riêng biệt:
  \begin{itemize}[nosep]
    \item \textbf{App Repo}: Chứa mã nguồn microservices, Dockerfile, Jenkinsfile. CI pipeline được kích hoạt khi có push vào repo này.
    \item \textbf{GitOps Repo}: Chứa các manifest triển khai (Kustomize/Helm). CI pipeline tự động cập nhật image tag vào repo này sau khi build thành công. Argo CD theo dõi repo này và đồng bộ lên cluster.
  \end{itemize}
\end{enumerate}

Mô hình dual repo mang lại lợi ích phân tách quyền: developer chỉ cần quyền trên App Repo, không cần truy cập GitOps Repo (vốn do CI pipeline tự động cập nhật). Điều này tăng cường bảo mật và giảm rủi ro.

\subsection{Nền tảng Kubernetes và Amazon EKS}

\subsubsection{Kubernetes}

Kubernetes (thường viết tắt là K8s) là một nền tảng mã nguồn mở để tự động hóa việc triển khai, scale, và quản lý các ứng dụng containerized [8]. Kubernetes được Google thiết kế dựa trên kinh nghiệm vận hành hệ thống Borg và được trao cho Cloud Native Computing Foundation (CNCF) quản lý từ năm 2015 [13].

Các thành phần chính của Kubernetes bao gồm:

\begin{itemize}[leftmargin=*,nosep]
  \item \textbf{Control Plane}: Bao gồm kube-apiserver, etcd (database phân tán), kube-controller-manager, và kube-scheduler. Quản lý trạng thái toàn cluster và ra quyết định (ví dụ: schedule pod lên node nào).
  \item \textbf{Worker Nodes}: Máy chủ chạy workload thực tế. Mỗi node chạy kubelet (agent giao tiếp với control plane), container runtime (Docker/containerd), và kube-proxy (quản lý network rules).
  \item \textbf{Pod}: Đơn vị triển khai nhỏ nhất trong Kubernetes, có thể chứa một hoặc nhiều container chia sẻ cùng network namespace và storage.
  \item \textbf{Service}: Một trừu tượng định nghĩa tập hợp các pod và policy để truy cập chúng, cung cấp IP ổn định và load balancing.
  \item \textbf{Deployment}: Quản lý việc triển khai và cập nhật pod, hỗ trợ rolling update và rollback.
  \item \textbf{Namespace}: Cơ chế phân vùng logic trong cluster, cho phép chia sẻ một cluster vật lý giữa nhiều team hoặc môi trường.
\end{itemize}

\subsubsection{Amazon EKS}

Amazon Elastic Kubernetes Service (EKS) là dịch vụ Kubernetes được quản lý bởi AWS, cung cấp control plane Kubernetes có tính sẵn sàng cao (highly available) trên nhiều Availability Zone [6]. Lợi ích của EKS so với tự quản lý Kubernetes:

\begin{itemize}[leftmargin=*,nosep]
  \item \textbf{Control plane được quản lý}: AWS đảm nhận việc vận hành, scale, và cập nhật control plane. Người dùng chỉ cần quản lý worker nodes.
  \item \textbf{Tích hợp sâu với hệ sinh thái AWS}: IAM Roles for Service Accounts (IRSA) cho phép gán IAM role trực tiếp cho Kubernetes service account; AWS Load Balancer Controller tự động tạo ALB/NLB; VPC-CNI cung cấp IP riêng cho mỗi pod.
  \item \textbf{Bảo mật}: Control plane chạy trong VPC riêng của AWS, tích hợp KMS để mã hóa secrets, hỗ trợ private endpoint.
  \item \textbf{Managed Node Groups}: AWS tự động quản lý việc provisioning, scaling, và cập nhật AMI cho worker nodes.
\end{itemize}

\subsection{Infrastructure as Code và Terraform}

Infrastructure as Code (IaC) là phương pháp quản lý hạ tầng CNTT thông qua các file cấu hình (code) thay vì thao tác thủ công qua giao diện (Console/CLI) [14]. Các nguyên tắc của IaC: idempotency (có thể chạy lại nhiều lần mà không thay đổi kết quả), declarative (khai báo trạng thái mong muốn, công cụ tự tính toán cách đạt được), và version control (lưu trong Git).

Terraform, phát triển bởi HashiCorp, là công cụ IaC phổ biến nhất hiện nay [7]. Terraform sử dụng ngôn ngữ HCL (HashiCorp Configuration Language) để định nghĩa tài nguyên, state file để theo dõi trạng thái thực tế, và provider để tương tác với các nền tảng khác nhau (AWS, Azure, GCP, Kubernetes, Helm).

Trong đồ án này, Terraform được sử dụng để quản lý toàn bộ hạ tầng AWS và Kubernetes addons. Mã Terraform được tổ chức thành các module độc lập (\texttt{vpc/}, \texttt{eks/}, \texttt{ecr/}) trong thư mục \texttt{03-infra/terraform/modules/}, mỗi module có biến đầu vào (variables.tf), đầu ra (outputs.tf), và khai báo provider (versions.tf) riêng biệt -- tuân theo các best practice về IaC của Terraform Registry:

\begin{itemize}[leftmargin=*,nosep]
  \item \textbf{Terraform Backend}: State được lưu trên S3 (có versioning), lock qua DynamoDB, cấu hình qua file \texttt{backend.hcl}.
  \item \textbf{VPC Module}: Sử dụng \texttt{terraform-aws-modules/vpc/aws} để tạo VPC, 6 subnets (3 public + 3 private), route tables, NAT Gateway.
  \item \textbf{EKS Module}: Sử dụng \texttt{terraform-aws-modules/eks/aws} để tạo EKS cluster, 3 managed node groups (2 t3.medium + 1 t3.large), VPC-CNI addon, IRSA, và các Kubernetes addons qua Helm provider.
  \item \textbf{ECR Module}: Module tự xây dựng (\texttt{modules/ecr/}) để tạo 5 ECR repositories cho từng microservice, cấu hình image scanning on push và lifecycle policy.
  \item Cấu trúc đa môi trường: \texttt{envs/dev/}, \texttt{envs/staging/}, \texttt{envs/prod/} -- mỗi môi trường có \texttt{terraform.tfvars} riêng.
\end{itemize}

\subsection{Các công cụ khác}

\subsubsection{Jenkins}

Jenkins là một automation server mã nguồn mở, được viết bằng Java, là nền tảng CI/CD phổ biến nhất trong giới doanh nghiệp [12]. Jenkins hỗ trợ kiến trúc plugin với hơn 1.800 plugin, cho phép tích hợp với hầu hết các công cụ trong hệ sinh thái DevOps. Trong đồ án này, Jenkins đảm nhiệm vai trò CI server:

\begin{itemize}[leftmargin=*,nosep]
  \item Lắng nghe webhook từ GitHub khi có push vào App Repo.
  \item Thực thi pipeline được định nghĩa trong Jenkinsfile (Pipeline as Code).
  \item Pipeline gồm các stage: Checkout → SonarQube Analysis → OWASP Dependency Check → Docker Build → Trivy Image Scan → Push ECR → Update GitOps Repo.
  \item Sử dụng Shared Libraries để đóng gói các hàm tái sử dụng: \texttt{awsSetup}, \texttt{buildService}, \texttt{sonarQubeAnalysis}, \texttt{owaspDependencyCheck}, \texttt{dockerBuildAndPush}, \texttt{updateGitOpsRepo} -- giúp Jenkinsfile chính chỉ còn 120 dòng so với hàng trăm dòng nếu viết trực tiếp.
  \item Docker images được build bằng Kaniko trong Kubernetes pods (không cần Docker daemon) để tăng tính bảo mật.
\end{itemize}

\subsubsection{Kustomize}

Kustomize là một công cụ quản lý cấu hình Kubernetes, được tích hợp sẵn trong kubectl từ phiên bản 1.14. Không giống như Helm (dùng template), Kustomize sử dụng phương pháp patch và overlay -- bắt đầu từ một bộ manifest cơ sở (base) và áp dụng các lớp tùy chỉnh (overlay) cho từng môi trường.

Cấu trúc Kustomize trong đồ án:

\begin{verbatim}
gitops-repo/apps/retail-store/
 -- base/
|    -- kustomization.yaml
|    -- ui-deployment.yaml
|    -- cart-deployment.yaml
|    -- orders-deployment.yaml
|    -- catalog-deployment.yaml
|    -- checkout-deployment.yaml
 -- overlays/
     -- dev/
    |    -- kustomization.yaml
    |    -- patches/
     -- staging/
      -- prod/
\end{verbatim}

\subsubsection{Helm}

Bên cạnh Kustomize, đồ án còn sử dụng Helm dưới dạng umbrella chart để cung cấp một cách tiếp cận khác cho việc quản lý cấu hình triển khai. Helm chart của ứng dụng retail-store được lưu trong GitOps repo tại \texttt{apps/retail-store/} với cấu trúc:

\begin{verbatim}
apps/retail-store/
 -- Chart.yaml
 -- values.yaml        # Default values
 -- templates/all.yaml # Định nghĩa Deployment, Service cho 5 service
\end{verbatim}

Argo CD được cấu hình với \texttt{--enable-helm} flag, cho phép sử dụng cả Kustomize và Helm trong cùng một Application. Tuy nhiên, trong bản demo hiện tại, Kustomize được sử dụng làm phương thức triển khai chính do tính linh hoạt và tích hợp sẵn trong kubectl.

\subsection{Policy as Code với Kyverno}

Policy as Code (PaC) là phương pháp quản lý và thực thi chính sách bảo mật, compliance thông qua code -- tương tự như IaC nhưng áp dụng cho các quy tắc vận hành và bảo mật [16]. Trong môi trường Kubernetes, Kyverno là công cụ PaC hàng đầu với các đặc điểm:

\begin{itemize}[leftmargin=*,nosep]
  \item \textbf{Kubernetes-native}: Chính sách được viết bằng YAML thuần túy, không cần học ngôn ngữ mới như Rego (OPA/Gatekeeper).
  \item \textbf{Ba chế độ hoạt động}:
  \begin{itemize}[nosep]
    \item \textit{Validate}: Chặn các resource không tuân thủ chính sách.
    \item \textit{Mutate}: Tự động sửa đổi resource (vd: thêm labels, security context).
    \item \textit{Generate}: Tự động sinh resource bổ sung (vd: tạo NetworkPolicy mặc định cho mỗi namespace).
  \end{itemize}
  \item \textbf{Audit mode}: Có thể chạy ở chế độ Audit (chỉ ghi log, không chặn) để kiểm tra tác động trước khi thực thi (Enforce).
  \item \textbf{Background scanning}: Tự động quét tất cả resource hiện có trong cluster.
\end{itemize}

Trong đồ án, 5 Kyverno ClusterPolicy được triển khai:
\begin{enumerate}[leftmargin=*,nosep]
  \item \texttt{require-pod-labels}: Yêu cầu tất cả pod phải có labels chuẩn \texttt{app.kubernetes.io/name} và \texttt{app.kubernetes.io/part-of}.
  \item \texttt{disallow-latest-tag}: Cấm sử dụng image tag \texttt{latest} trong production, buộc phải dùng tag cụ thể (git SHA).
  \item \texttt{require-probes}: Yêu cầu mọi Deployment phải khai báo liveness và readiness probe.
  \item \texttt{require-resource-limits}: Yêu cầu tất cả container phải có resource requests và limits.
  \item \texttt{disallow-root-containers}: Cấm chạy container với quyền root (UID 0), tăng cường container security.
\end{enumerate}

\subsection{Sample Application: Retail Store}

Đồ án sử dụng sample application từ AWS: \texttt{aws-containers/retail-store-sample-app} [15]. Đây là một ứng dụng bán lẻ trực tuyến với kiến trúc microservices, được AWS thiết kế để minh họa các best practice trong việc triển khai ứng dụng trên EKS.

\begin{table}[H]
\centering
\begin{tabular}{@{}p{2cm} p{2.5cm} p{4cm} p{5.5cm}@{}}
\toprule
\textbf{Service} & \textbf{Ngôn ngữ} & \textbf{Framework} & \textbf{Backend/Phụ thuộc} \\
\midrule
UI & Java 21 & Spring Boot + Thymeleaf & WebFlux API Gateway \\
Cart & Java 21 & Spring Boot & DynamoDB \\
Orders & Java 21 & Spring Boot & DynamoDB, SQS, SNS, Secrets Manager \\
Catalog & Go & Gin & MySQL / MariaDB \\
Checkout & TypeScript & NestJS 11 & Redis \\
\bottomrule
\end{tabular}
\caption{Các microservice trong Retail Store Sample App}
\label{tab:services}
\end{table}

Ứng dụng được chọn vì các lý do: đa dạng ngôn ngữ (Java, Go, TypeScript) giúp thể hiện tính linh hoạt của CI pipeline; kiến trúc microservices đầy đủ với database riêng cho từng service; tích hợp sẵn với các dịch vụ AWS (DynamoDB, SQS, SNS, Secrets Manager); có sẵn Dockerfile và Terraform module.

\newpage

% ════════════════════════════════════════════════════
%  CHƯƠNG 3: PHƯƠNG PHÁP VÀ THIẾT KẾ HỆ THỐNG
% ════════════════════════════════════════════════════
\section{Phương pháp và Thiết kế Hệ thống}

Chương này trình bày kiến trúc tổng thể của hệ thống, thiết kế chi tiết từng thành phần, và phương pháp triển khai.

\subsection{Kiến trúc tổng thể}

Hình \ref{fig:architecture} trình bày kiến trúc tổng thể của hệ thống DevSecOps + GitOps.

\begin{figure}[H]
\centering
\begin{verbatim}
  ------------------------------------------------------------- 
 |                    CI Pipeline (Jenkins)                     |
 |                                                             |
 |   -------     ----------     --------     --------         |
 |  |  Code |-->|SonarQube |-->| OWASP  |-->| Docker |        |
 |  | Check |   |   SAST   |   |   DC   |   | Build  |        |
 |   -------     ----------     --------     ---- ---         |
 |                                               v             |
 |                                     --------------------    |
 |  Developer                          |   Trivy Image Scan |   |
 |  --> Push code                       -------- -----------    |
 |       to App Repo                            v               |
 |                                      --------------------    |
 |                                     |   Push to ECR      |   |
 |                                      -------- -----------    |
  ----------------------------------------------+--------------- 
                                                v
  ------------------------------------------------------------- 
 |                    CD GitOps (Argo CD)                       |
 |                                                             |
 |   ------------------           -------------------------     |
 |  |  GitOps Repo      |  Sync   |  EKS Cluster             |    |
 |  |  (Kustomize)      |<------>|   ---------------------  |    |
 |  |                   |         |  |  Managed Node Group  | |    |
 |  |  base/            |         |  |   -----   -----     | |    |
 |  |  overlays/dev/    |         |  |  | Pod | | Pod |... | |    |
 |   ------------------          |  |   -----   -----     | |    |
 |                               |   ---------------------  |    |
 |                               |   ---------------------  |    |
 |                               |  | Load Balancer       | |    |
 |                               |  | Controller + ALB    | |    |
 |                               |   ---------------------  |    |
 |                                -------------------------     |
  ------------------------------------------------------------- 
\end{verbatim}
\caption{Kiến trúc tổng thể DevSecOps + GitOps}
\label{fig:architecture}
\end{figure}

\subsection{Cấu trúc Workspace đồ án}

Workspace được tổ chức thành 10 danh mục, mỗi danh mục đảm nhiệm một vai trò cụ thể:

\begin{table}[H]
\centering
\begin{tabular}{@{}p{3.5cm} p{11.5cm}@{}}
\toprule
\textbf{Thư mục} & \textbf{Mục đích} \\
\midrule
00-admin & Quản lý dự án: meeting notes, decision log, checklist. \\
01-docs & Tài liệu: báo cáo tiến độ, báo cáo tổng kết, slides bảo vệ, đề cương. \\
02-repos & Mã nguồn: App Repo (sample app) và GitOps Repo (Kustomize manifests). \\
03-infra & Infrastructure as Code: Terraform backend, môi trường dev/staging/prod. \\
04-platform & Kubernetes platform: cấu hình Argo CD, Helm values. \\
05-ci & CI/CD pipeline: Jenkinsfile, Shared Libraries (6 Groovy functions), security gates, Kaniko configs. \\
06-security & Security artifacts: Kyverno ClusterPolicies (5), Terrascan scan reports, Checkov setup. \\
07-observability & Monitoring: Grafana dashboard JSON, Prometheus alerting rules (10 rules). \\
08-evidence & Bằng chứng: screenshots, logs, kết quả kiểm thử từng phase (phase-2-infra, phase-3-ci, phase-4-gitops, phase-4-security, phase-4-observability). \\
09-scripts & Script hỗ trợ: setup toolchain, switch AWS profile. \\
\bottomrule
\end{tabular}
\caption{Cấu trúc workspace đồ án}
\label{tab:workspace}
\end{table}

\subsection{Thiết kế Hạ tầng AWS}

\subsubsection{Sơ đồ mạng VPC}

VPC được thiết kế theo nguyên tắc phân tách public/private subnet:

\begin{table}[H]
\centering
\begin{tabular}{@{}p{4cm} p{5cm} p{6cm}@{}}
\toprule
\textbf{Subnet} & \textbf{Số lượng/CIDR} & \textbf{Mục đích} \\
\midrule
Public Subnet & 3 (mỗi /24, 3 AZ) & Chứa Load Balancer (ALB/NLB), Internet-facing. Có route đến Internet Gateway. \\
\midrule
Private Subnet & 3 (mỗi /24, 3 AZ) & Chứa EKS worker nodes. Không có IP public trực tiếp. Truy cập internet qua NAT Gateway. \\
\bottomrule
\end{tabular}
\caption{Thiết kế subnet trong VPC}
\label{tab:subnet-design}
\end{table}

Bảng \ref{tab:vpc-config} trình bày các thông số cấu hình chính của VPC đã triển khai.

\begin{table}[H]
\centering
\begin{tabular}{@{}p{5cm} p{10cm}@{}}
\toprule
\textbf{Thông số} & \textbf{Giá trị} \\
\midrule
VPC ID & vpc-0d37b4e9b21e95421 \\
CIDR Block & 10.0.0.0/16 \\
Internet Gateway & igw-089e9cb13e3af83dd \\
NAT Gateway & 1 (single, dùng chung) \\
Public Route Table & rtb-0331cb9ad33461b8c \\
Private Route Table & rtb-0dac22b57694f1444 \\
EKS Cluster Tag & kubernetes.io/cluster/uit-devsecops-dev = shared \\
ELB Tag (public) & kubernetes.io/role/elb = 1 \\
ELB Tag (internal) & kubernetes.io/role/internal-elb = 1 \\
\bottomrule
\end{tabular}
\caption{Cấu hình VPC đã triển khai}
\label{tab:vpc-config}
\end{table}

\subsubsection{EKS Cluster}

Bảng \ref{tab:eks-config} trình bày cấu hình chi tiết của EKS cluster.

\begin{table}[H]
\centering
\begin{tabular}{@{}p{5cm} p{10cm}@{}}
\toprule
\textbf{Thông số} & \textbf{Giá trị} \\
\midrule
Tên cluster & uit-devsecops-dev \\
Phiên bản Kubernetes & v1.33.11 \\
Control plane endpoint & Public + Private \\
Bảo mật etcd & KMS encryption (key alias: eks/uit-devsecops-dev) \\
OIDC Provider & Đã tạo (hỗ trợ IRSA) \\
CloudWatch Logs & /aws/eks/uit-devsecops-dev/cluster \\
\midrule
\multicolumn{2}{c}{\textbf{Managed Node Groups}} \\
\midrule
Node Group 1 & managed-nodegroup-1 (AZ 1) \\
Node Group 2 & managed-nodegroup-2 (AZ 2) \\
Node Group 3 & managed-nodegroup-large (AZ 1, dành cho workload nặng) \\
Instance Type (small) & t3.medium (2 vCPU, 4GB RAM) \\
Instance Type (large) & t3.large (2 vCPU, 8GB RAM) \\
Capacity Type & ON\_DEMAND \\
AMI & Amazon Linux 2023 optimized for EKS \\
Scaling (min/desired/max) & 0/0/1 (desired thay đổi theo nhu cầu, scale về 0 để tiết kiệm) \\
Total nodes khi active & 3 nodes (2 t3.medium + 1 t3.large) \\
IAM Policy & AmazonEKSWorkerNodePolicy \\
 & AmazonEKS\_CNI\_Policy \\
 & AmazonEC2ContainerRegistryReadOnly \\
\bottomrule
\end{tabular}
\caption{Cấu hình EKS cluster}
\label{tab:eks-config}
\end{table}

\subsubsection{Kubernetes Addons}

Bảng \ref{tab:addons} liệt kê các addon đã được cài đặt trên cluster.

\begin{table}[H]
\centering
\begin{tabular}{@{}p{3cm} p{3cm} p{2.5cm} p{5.5cm}@{}}
\toprule
\textbf{Addon} & \textbf{Namespace} & \textbf{Replicas} & \textbf{Chức năng} \\
\midrule
VPC-CNI & managed & 2 (1/node) & Quản lý IP pod qua ENI \\
CoreDNS & kube-system & 2 & DNS nội bộ cluster \\
kube-proxy & kube-system & 2 & Network proxy, Service \\
AWS LB Controller & kube-system & 2 (HA) & Tự động tạo ALB/NLB \\
Cert Manager & cert-manager & 3 & Cấp phát TLS certificate \\
OTEL Operator & opentelemetry-operator-system & 1 & OpenTelemetry (đã disable — không cần thiết cho demo) \\
\bottomrule
\end{tabular}
\caption{Danh sách Kubernetes addons đã cài đặt}
\label{tab:addons}
\end{table}

\subsection{Thiết kế Security Pipeline}

Hình \ref{fig:security-pipeline} mô tả luồng security checks trong CI pipeline.

\begin{figure}[H]
\centering
\begin{verbatim}
    ----------------------------------------------------- 
   |                    CI Pipeline                       |
   |                                                     |
   |  Push Code                                           |
   |     |                                                |
   |     v                                                |
   |   ----------------------                             |
   |  |  Checkout + Build     |                           |
   |   ---------- -----------                             |
   |             v                                        |
   |   ----------------------                             |
   |  | SECURITY GATE 1       | SAST                      |
   |  | SonarQube Analysis    | --> PASS/FAIL              |
   |   ---------- -----------     |                        |
   |             | PASS            | FAIL → STOP            |
   |             v                                        |
   |   ----------------------                             |
   |  | SECURITY GATE 2       | SCA                       |
   |  | OWASP Dependency Check| --> PASS/FAIL              |
   |   ---------- -----------     |                        |
   |             | PASS            | FAIL → STOP            |
   |             v                                        |
   |   ----------------------                             |
   |  | Docker Build          |                           |
   |   ---------- -----------                             |
   |             v                                        |
   |   ----------------------                             |
   |  | SECURITY GATE 3       | Container Scan             |
   |  | Trivy Image Scan      | --> PASS/FAIL              |
   |   ---------- -----------     |                        |
   |             | PASS            | FAIL → STOP            |
   |             v                                        |
   |   ----------------------                             |
   |  | Push to ECR           | <-- Chỉ image sạch         |
   |   ---------- -----------     được push                |
   |             v                                        |
   |   ----------------------                             |
   |  | Update GitOps Repo    | Cập nhật image tag mới     |
   |   ----------------------                             |
   |                                                     |
   |  Chỉ artifact vượt qua CẢ BA security gate          |
   |  mới được push lên registry                         |
    ----------------------------------------------------- 
\end{verbatim}
\caption{Thiết kế CI pipeline với 3 security gates}
\label{fig:security-pipeline}
\end{figure}

Mỗi security gate có ngưỡng thất bại (failure threshold) được định nghĩa:

\begin{table}[H]
\centering
\begin{tabular}{@{}p{3.5cm} p{4cm} p{7.5cm}@{}}
\toprule
\textbf{Security Gate} & \textbf{Ngưỡng} & \textbf{Hành động khi không đạt} \\
\midrule
SonarQube SAST & Quality Gate FAIL & Pipeline dừng, gửi thông báo đến developer \\
OWASP DC & CVSS ≥ 7.0 (HIGH) & Pipeline dừng, báo cáo CVE được lưu \\
Trivy Scan & CRITICAL CVE & Pipeline dừng, image KHÔNG được push \\
\bottomrule
\end{tabular}
\caption{Ngưỡng cho từng security gate}
\label{tab:security-gates}
\end{table}

\subsection{Thiết kế GitOps Pipeline}

\subsubsection{Mô hình 2-Repository}

Hình \ref{fig:2repo} minh họa mô hình 2-repository được triển khai trong đồ án.

\begin{figure}[H]
\centering
\begin{verbatim}
    -----------------           ------------------ 
   |   App Repo       |         |  GitOps Repo      |
   |                  |         |                   |
   |  src/ui/         |         |  apps/store/      |
   |  src/cart/       |         |   -- base/        |
   |  src/orders/     |         |   -- overlays/    |
   |  src/catalog/    |         |       -- dev/     |
   |  src/checkout/   |         |       -- staging/ |
   |  Dockerfile      |         |       -- prod/    |
   |  Jenkinsfile     |         |                   |
   |                  |         |                   |
    -------- --------           ---------^--------- 
            |                            |
            | 1. Push code               |
            v                            |
    -----------------                    |
   |    Jenkins CI     |                  |
   |                  | 3. git commit    |
   |  Build → Scan    |---- Update ------ 
   |  → Push ECR      |   image tag
    ----------------- 
\end{verbatim}
\caption{Mô hình 2-repository GitOps}
\label{fig:2repo}
\end{figure}

\subsubsection{Kustomize Base \& Overlay}

Cấu trúc Kustomize theo mô hình base + overlay:

\begin{itemize}[leftmargin=*,nosep]
  \item \textbf{Base}: Định nghĩa cấu hình chung cho tất cả môi trường: Deployment, Service, ConfigMap cơ bản.
  \item \textbf{Overlay dev}: Ghi đè các giá trị cho môi trường dev: replica count = 1, resource limits thấp, image tag từ CI.
  \item \textbf{Overlay staging}: (chuẩn bị) replica count = 2, resource limits trung bình.
  \item \textbf{Overlay prod}: (chuẩn bị) replica count = 3+, HPA enabled, stricter resource limits.
\end{itemize}

\subsubsection{Argo CD Application}

Mỗi service được định nghĩa là một Argo CD Application với cấu hình:

\begin{verbatim}
apiVersion: argoproj.io/v1alpha1
kind: Application
metadata:
  name: retail-store-dev
  namespace: argocd
spec:
  project: retail-store
  source:
    repoURL: https://github.com/<org>/gitops-repo.git
    targetRevision: HEAD
    path: apps/retail-store/overlays/dev
  destination:
    server: https://kubernetes.default.svc
    namespace: retail-store-dev
  syncPolicy:
    automated:
      prune: true
      selfHeal: true
    syncOptions:
      - CreateNamespace=true
\end{verbatim}

\subsection{Thiết kế ECR Repositories}

Năm ECR repository được tạo cho từng microservice, mỗi repository có cấu hình:

\begin{itemize}[leftmargin=*,nosep]
  \item \textbf{Image scanning on push}: Mỗi lần push image, ECR tự động scan CVE bằng Clair (engine tích hợp sẵn của AWS).
  \item \textbf{Lifecycle policy}: Tự động xóa các image cũ, chỉ giữ lại 10 image mới nhất để tiết kiệm chi phí lưu trữ.
  \item \textbf{Tag mutability}: Cho phép ghi đè tag (MUTABLE) trong giai đoạn dev.
\end{itemize}

\newpage

% ════════════════════════════════════════════════════
%  CHƯƠNG 4: KẾT QUẢ TRIỂN KHAI VÀ KIỂM THỬ
% ════════════════════════════════════════════════════
\section{Kết quả Triển khai và Kiểm thử}

Chương này trình bày kết quả triển khai thực tế trên AWS và các bước kiểm thử đã thực hiện.

\subsection{Triển khai hạ tầng}

\subsubsection{Kết quả Terraform Apply}

Toàn bộ hạ tầng được provision bằng lệnh \texttt{terraform apply} trong một lần duy nhất. Bảng \ref{tab:tf-results} tổng hợp kết quả.

\begin{table}[H]
\centering
\begin{tabular}{@{}p{5cm} p{10cm}@{}}
\toprule
\textbf{Thông số} & \textbf{Giá trị} \\
\midrule
Thời gian provision & ~13 phút \\
Số resource được tạo & 101 resources (bao gồm VPC + EKS + Addons + ECR) \\
Số resource thay đổi & 0 \\
Số resource bị hủy & 0 \\
State file & Lưu trên S3: \texttt{s3://uit-devsecops-gitops-dev-0a502713-tfstate/envs/dev/terraform.tfstate} \\
Lock & DynamoDB: \texttt{uit-devsecops-gitops-dev-0a502713-tflock} \\
\bottomrule
\end{tabular}
\caption{Kết quả Terraform Apply}
\label{tab:tf-results}
\end{table}

Quá trình provision diễn ra tuần tự theo dependency graph của Terraform: VPC được tạo đầu tiên, sau đó đến EKS cluster và VPC-CNI addon, tiếp đến managed node groups, và cuối cùng là Helm releases (Load Balancer Controller, Cert Manager) và ECR repositories.

\evidence{
\textbf{Evidence E01 -- Terraform Apply thành công}

\smallskip
\textbf{Cách lấy:} Mở terminal, chạy lại lệnh \texttt{terraform apply} (hoặc dùng output đã lưu). Chụp màn hình dòng:
\begin{verbatim}
Apply complete! Resources: 101 added, 0 changed, 0 destroyed.
\end{verbatim}
\textbf{Ngoài ra chụp thêm:}
\begin{itemize}[leftmargin=*,nosep]
  \item AWS Console \textrightarrow\ EKS \textrightarrow\ Clusters \textrightarrow\ \texttt{uit-devsecops-dev} (thể hiện status Active).
  \item AWS Console \textrightarrow\ VPC \textrightarrow\ Your VPCs (thể hiện VPC \texttt{uit-devsecops-dev} với 6 subnets).
  \item AWS Console \textrightarrow\ EC2 \textrightarrow\ Instances (thể hiện 3 instances: 2 t3.medium + 1 t3.large).
  \item AWS Console \textrightarrow\ ECR \textrightarrow\ Repositories (thể hiện 5 repos).
\end{itemize}
\textbf{Lưu vào:} \texttt{08-evidence/phase-2-infra/}
}

\subsubsection{Kết quả kiểm tra cluster}

Sau khi provision hoàn tất, các bước kiểm tra kết nối và trạng thái cluster được thực hiện.

\begin{verbatim}
$ kubectl cluster-info
Kubernetes control plane is running at
  https://F358B0B27A319352351851239024C3F2.gr7.ap-southeast-1.eks.amazonaws.com
CoreDNS is running at
  https://.../api/v1/namespaces/kube-system/services/kube-dns:dns/proxy
\end{verbatim}

\begin{verbatim}
$ kubectl get nodes -o wide
NAME                    STATUS  ROLES   VERSION              INTERNAL-IP
ip-10-0-10-249...       Ready   <none>  v1.33.11-eks-4136f65 10.0.10.249
ip-10-0-11-155...       Ready   <none>  v1.33.11-eks-4136f65 10.0.11.155
ip-10-0-12-45...        Ready   <none>  v1.33.11-eks-4136f65 10.0.12.45
\end{verbatim}

Cả ba worker node (2 t3.medium + 1 t3.large) đều ở trạng thái Ready, chạy Amazon Linux 2023, container runtime containerd, với địa chỉ IP nằm trong dải private subnet (10.0.0.0/16). Node t3.large được sử dụng cho các workload nặng hơn như Jenkins và Prometheus/Grafana.

\begin{verbatim}
$ kubectl get pods -A
NAMESPACE       NAME                               READY   STATUS    RESTARTS
argocd          argocd-application-controller-0     1/1     Running   0
argocd          argocd-applicationset-...           1/1     Running   0
argocd          argocd-dex-server-...               1/1     Running   0
argocd          argocd-notifications-...            1/1     Running   0
argocd          argocd-redis-...                    1/1     Running   0
argocd          argocd-repo-server-...              1/1     Running   0
argocd          argocd-server-...                   1/1     Running   0
cert-manager    cert-manager-...                    1/1     Running   0
cert-manager    cert-manager-cainjector-...         1/1     Running   0
cert-manager    cert-manager-webhook-...            1/1     Running   0
jenkins         jenkins-0                           2/2     Running   0
kube-system     aws-load-balancer-controller-...    1/1     Running   0
kube-system     aws-load-balancer-controller-...    1/1     Running   0
kube-system     aws-node-s9glr                      2/2     Running   0
kube-system     aws-node-whlpn                      2/2     Running   0
kube-system     aws-node-xyz12                      2/2     Running   0
kube-system     coredns-...                         1/1     Running   0
kube-system     coredns-...                         1/1     Running   0
kube-system     kube-proxy-...                      1/1     Running   0
kube-system     kube-proxy-...                      1/1     Running   0
kube-system     kube-proxy-...                      1/1     Running   0
monitoring      alertmanager-...                    1/1     Running   0
monitoring      kube-prometheus-stack-grafana-...   3/3     Running   0
monitoring      kube-prometheus-stack-operator-...  1/1     Running   0
monitoring      kube-prometheus-stack-prometheus-...2/2     Running   0
monitoring      kube-state-metrics-...              1/1     Running   0
monitoring      prometheus-node-exporter-...        1/1     Running   0
monitoring      prometheus-node-exporter-...        1/1     Running   0
monitoring      prometheus-node-exporter-...        1/1     Running   0
retail-store-dev ui-...                             1/1     Running   0
retail-store-dev ui-...                             1/1     Running   0
retail-store-dev cart-...                           1/1     Running   0
retail-store-dev orders-...                         1/1     Running   0
retail-store-dev catalog-...                        1/1     Running   0
retail-store-dev checkout-...                       1/1     Running   0
\end{verbatim}

Tổng cộng 41 pods, trong đó 26 pod hệ thống (Argo CD, Jenkins, Prometheus, Grafana, Cert Manager, Kyverno, SonarQube, CoreDNS, VPC-CNI, kube-proxy, LB Controller) và 7 pod ứng dụng (UI ×2, Cart, Orders, Catalog, Catalog-DB, Checkout) đều ở trạng thái Running. Toàn bộ pod hệ thống hoạt động ổn định, không có CrashLoopBackOff.

\evidence{
\textbf{Evidence E02 -- Kết nối và kiểm tra EKS cluster}

\smallskip
\textbf{Cách lấy:} Mở terminal, chạy và chụp màn hình các lệnh:
\begin{itemize}[leftmargin=*,nosep]
  \item \texttt{kubectl cluster-info} \textrightarrow\ thể hiện control plane endpoint.
  \item \texttt{kubectl get nodes -o wide} \textrightarrow\ thể hiện 3 nodes Ready (2 t3.medium + 1 t3.large).
  \item \texttt{kubectl get pods -A} \textrightarrow\ thể hiện 41 pods Running.
  \item \texttt{kubectl get svc -A} \textrightarrow\ thể hiện các service đã tạo.
\end{itemize}
\textbf{Lưu vào:} \texttt{08-evidence/phase-2-infra/}
}

\subsection{Kiểm tra ECR Repositories}

Năm ECR repository đã được tạo thành công, xác nhận qua AWS Console và Terraform output.

\begin{table}[H]
\centering
\begin{tabular}{@{}p{9.5cm} p{5.5cm}@{}}
\toprule
\textbf{Service} & \textbf{Repository} \\
\midrule
UI & uit-devsecops-dev-ui \\
Cart & uit-devsecops-dev-cart \\
Orders & uit-devsecops-dev-orders \\
Catalog & uit-devsecops-dev-catalog \\
Checkout & uit-devsecops-dev-checkout \\
\bottomrule
\end{tabular}
\caption{ECR repositories đã tạo}
\label{tab:ecr-repos}
\end{table}

Mỗi repository được cấu hình Image Scanning on Push và Lifecycle Policy giữ 10 image mới nhất.

\subsection{Kiểm tra trên AWS Console}

\begin{itemize}[leftmargin=*,nosep]
  \item \textbf{EKS Dashboard}: Cluster \texttt{uit-devsecops-dev} trạng thái Active, 3 managed node groups hiển thị Active.
  \item \textbf{EC2 Dashboard}: 3 instances (2 t3.medium + 1 t3.large) chạy trong private subnet.
  \item \textbf{VPC Dashboard}: VPC đầy đủ 6 subnets, Internet Gateway, NAT Gateway, route tables.
  \item \textbf{CloudWatch}: Log group \texttt{/aws/eks/uit-devsecops-dev/cluster} đang nhận log.
  \item \textbf{ECR}: 5 repositories đã tạo, sẵn sàng nhận Docker images.
  \item \textbf{IAM}: Cluster role, node role, addon roles (LB Controller, Cert Manager) đầy đủ.
\end{itemize}

\evidence{
\textbf{Evidence E03 -- AWS Console Verification (ECR, EKS, VPC, EC2, IAM)}

\smallskip
\textbf{Cách lấy:} Vào 5 AWS Console dưới đây, chụp màn hình:
\begin{itemize}[leftmargin=*,nosep]
  \item \textbf{ECR}: AWS Console \textrightarrow\ ECR \textrightarrow\ Repositories \textrightarrow\ chụp danh sách 5 repos (\texttt{uit-devsecops-dev-*}).
  \item \textbf{EKS}: AWS Console \textrightarrow\ EKS \textrightarrow\ Clusters \textrightarrow\ \texttt{uit-devsecops-dev} \textrightarrow\ tab Overview + tab Compute (thấy 3 node groups).
  \item \textbf{VPC}: AWS Console \textrightarrow\ VPC \textrightarrow\ Your VPCs \textrightarrow\ chụp VPC \texttt{uit-devsecops-dev} + tab Subnets.
  \item \textbf{EC2}: AWS Console \textrightarrow\ EC2 \textrightarrow\ Instances \textrightarrow\ chụp 3 instances (2 t3.medium + 1 t3.large).
  \item \textbf{IAM}: AWS Console \textrightarrow\ IAM \textrightarrow\ Roles \textrightarrow\ tìm \texttt{uit-devsecops-dev*} \textrightarrow\ chụp danh sách role.
\end{itemize}
\textbf{Lưu vào:} \texttt{08-evidence/phase-2-infra/}
}

\subsection{Triển khai CI Pipeline}

Jenkins đã được triển khai trên EKS cluster dưới dạng StatefulSet trong namespace \texttt{jenkins} với 1 pod (2/2 containers: Jenkins master + Jenkins agent). Jenkins được cấu hình với:

\begin{itemize}[leftmargin=*,nosep]
  \item \textbf{Shared Library}: Thư viện dùng chung \texttt{devsecops-shared-library} được load từ GitHub, chứa 6 hàm Groovy trong \texttt{05-ci/jenkins-shared-library/vars/}.
  \item \textbf{AWS Credentials}: Sử dụng CloudBees AWS Credentials plugin, cấu hình qua IAM role (IRSA) thay vì hard-code Access Key/Secret Key.
  \item \textbf{Kaniko Build}: Docker images được build bằng Kaniko chạy trong Kubernetes pod (pod template agent). Kaniko không yêu cầu Docker daemon, tăng tính bảo mật và phù hợp với môi trường Kubernetes.
  \item \textbf{Pipeline as Code}: Jenkinsfile được lưu trong App Repo, Jenkins tự động fetch và thực thi khi có webhook từ GitHub.
\end{itemize}

CI pipeline đã được thực thi thành công với Kaniko builds cho cả 5 microservice, vượt qua 3 security gate (SonarQube SAST, OWASP DC, Trivy), và tự động push Docker images lên ECR với tag dựa trên git commit SHA. Sau khi push ECR thành công, pipeline tự động cập nhật image tag trong GitOps repo thông qua script \texttt{updateGitOpsRepo.groovy}, kích hoạt Argo CD tự động sync lên EKS cluster.

Kết quả các lần chạy CI pipeline được lưu trữ trong Jenkins build history, bao gồm artifact (Trivy scan reports, OWASP dependency check reports) và logs của từng stage.

\subsubsection{Cấu trúc file CI (đã triển khai)}

Pipeline được triển khai dưới dạng Declarative Pipeline của Jenkins, quản lý dưới dạng Pipeline as Code (Jenkinsfile). Pipeline sử dụng Jenkins Shared Libraries để đóng gói các hàm tái sử dụng, giúp Jenkinsfile chính ngắn gọn (136 dòng) và dễ bảo trì. Sáu hàm Groovy được định nghĩa trong thư mục \texttt{05-ci/jenkins-shared-library/vars/}:

\begin{itemize}[leftmargin=*,nosep]
  \item \texttt{awsSetup.groovy}: Thiết lập AWS credentials từ CloudBees AWS Credentials plugin (không hard-code Access Key).
  \item \texttt{buildService.groovy}: Build source code tự động phát hiện ngôn ngữ (Maven cho Java, Yarn cho Node.js, Go build).
  \item \texttt{sonarQubeAnalysis.groovy}: Gọi SonarQube Scanner API, kiểm tra Quality Gate status.
  \item \texttt{owaspDependencyCheck.groovy}: Chạy OWASP Dependency Check, phân tích CVSS score.
  \item \texttt{dockerBuildAndPush.groovy}: Build Docker image bằng Kaniko trong Kubernetes pod (không cần Docker daemon), quét bằng Trivy, push lên ECR.
  \item \texttt{updateGitOpsRepo.groovy}: Clone GitOps repo, cập nhật image tag trong các file Kustomize/Helm, commit và push.
\end{itemize}

Dưới đây là danh sách đầy đủ các stage trong Jenkinsfile chính:

\begin{table}[H]
\centering
\begin{tabular}{@{}p{1cm} p{4.5cm} p{8.5cm}@{}}
\toprule
\textbf{\#} & \textbf{Stage} & \textbf{Mô tả} \\
\midrule
1 & Checkout      & Clone mã nguồn từ App Repo, xác định branch và image tag. \\
2 & AWS Setup     & Cấu hình AWS credentials qua CloudBees plugin (dùng IAM role, không hard-code). \\
3 & Build         & Build song song 5 service (parallel): Maven (ui, cart, orders), Go (catalog), Yarn (checkout). \\
4 & SonarQube SAST & SECURITY GATE 1 --- Phân tích mã nguồn tĩnh, chặn nếu Quality Gate FAIL. \\
5 & OWASP DC      & SECURITY GATE 2 --- Quét dependency, chặn nếu có CVE với CVSS ≥ 7.0. \\
6 & Docker Build, Scan \& Push & Build Docker image bằng Kaniko (không cần Docker daemon), quét Trivy (SECURITY GATE 3), push ECR. \\
7 & Update GitOps & Cập nhật image tag trong GitOps repo, commit và push để Argo CD sync. \\
\bottomrule
\end{tabular}
\caption{Các stage trong Jenkins CI Pipeline}
\label{tab:jenkins-stages}
\end{table}

\subsubsection{Security Gates}

Ba security gate được tích hợp vào pipeline theo nguyên tắc ``fail fast'' -- pipeline dừng ngay khi phát hiện lỗ hổng, không tiếp tục các bước sau:

\textbf{SonarQube SAST} phân tích toàn bộ mã nguồn với các rule về security, reliability, maintainability. Quality Gate được cấu hình với các điều kiện: security\_rating không được tệ hơn C, reliability\_rating tối thiểu B, new\_coverage tối thiểu 50\%, new\_duplicated\_lines\_density tối đa 10\%.

\textbf{OWASP Dependency Check} quét tất cả file cấu hình dependency (pom.xml, package.json, go.mod) và đối chiếu với NVD database. Pipeline bị chặn nếu có bất kỳ dependency nào có CVSS ≥ 7.0 (HIGH hoặc CRITICAL).

\textbf{Trivy Image Scan} quét Docker image đã build, kiểm tra OS packages và application dependencies. Chỉ chặn pipeline khi có CRITICAL CVE (không chặn với HIGH, chỉ cảnh báo).

\subsubsection{Quản lý Image Tag}

Mỗi lần pipeline chạy, image được gán tag dựa trên git commit SHA (7 ký tự đầu). Đồng thời tag \texttt{latest} cũng được cập nhật. Tag được lưu trong environment variable \texttt{IMAGE\_TAG} và có thể override qua pipeline parameter.

\subsubsection{Cập nhật GitOps Repo}

Sau khi push ECR thành công, hàm \texttt{updateGitOpsRepo.groovy} (trong Shared Library) clone GitOps repo, tìm các file kustomization.yaml trong thư mục \texttt{overlays/dev/} và cập nhật trường \texttt{newTag} cho từng image tương ứng. Sau đó commit với message ``ci: update image tags to <sha> [skip ci]'' và push. Argo CD phát hiện commit mới và tự động sync lên EKS.

\subsubsection{Cấu trúc file CI}

\begin{verbatim}
05-ci/
  jenkins/
    Jenkinsfile                     # Pipeline chinh (declarative, 7 stage, @Library)
    update-gitops-repo.py           # Script Python cap nhat GitOps repo (legacy)
    README.md
  jenkins-shared-library/
    vars/
      awsSetup.groovy               # Cấu hình AWS credentials (IRSA)
      buildService.groovy           # Build đa ngôn ngữ (Maven/Go/Yarn)
      sonarQubeAnalysis.groovy      # SonarQube SAST + Quality Gate
      owaspDependencyCheck.groovy   # OWASP Dependency Check
      dockerBuildAndPush.groovy     # Kaniko build + Trivy scan + Push ECR
      updateGitOpsRepo.groovy       # Cập nhật GitOps repo + commit push
  security-gates/
    sonarqube-quality-gate.yml      # Quality Gate conditions
    owasp-threshold.yml             # CVSS threshold >= 7.0
    trivy-policy.yml                # Severity: CRITICAL + HIGH
  README.md                         # Setup guide
\end{verbatim}

\evidence{
\textbf{Evidence E04 -- CI Pipeline (Jenkinsfile + Security Gates)}

\smallskip
\textbf{Cách lấy:} Chụp màn hình các file sau trong IDE (VS Code):
\begin{itemize}[leftmargin=*,nosep]
  \item \texttt{05-ci/jenkins/Jenkinsfile} \textrightarrow\ chụp toàn bộ file (thể hiện 7 stage, @Library import).
  \item \texttt{05-ci/jenkins-shared-library/vars/} \textrightarrow\ chụp 6 file Groovy (awsSetup, buildService, sonarQubeAnalysis, owaspDependencyCheck, dockerBuildAndPush, updateGitOpsRepo).
  \item \texttt{05-ci/security-gates/sonarqube-quality-gate.yml} \textrightarrow\ cấu hình Quality Gate.
  \item \texttt{05-ci/security-gates/owasp-threshold.yml} \textrightarrow\ threshold OWASP.
  \item \texttt{05-ci/security-gates/trivy-policy.yml} \textrightarrow\ policy Trivy.
\end{itemize}
\textbf{Lưu vào:} \texttt{08-evidence/phase-3-ci/}
}

\subsection{Triển khai CD GitOps với Argo CD}

\subsubsection{Kustomize Manifests}

Manifest triển khai cho 5 microservice được tổ chức theo mô hình Kustomize base + overlay. Base định nghĩa Deployment và Service chung cho tất cả môi trường, overlay dev ghi đè các giá trị riêng như replica count, resource limits và namespace.

Bảng \ref{tab:kustomize-base} liệt kê các file trong base.

\begin{table}[H]
\centering
\begin{tabular}{@{}p{6cm} p{8cm}@{}}
\toprule
\textbf{File} & \textbf{Nội dung} \\
\midrule
base/kustomization.yaml & Danh sách resources, ánh xạ image name → ECR URL, common labels \\
base/deployment.yaml & Deployment cho cả 5 service (UI, Cart, Orders, Catalog, Checkout) \\
base/service.yaml & Service (ClusterIP) cho cả 5 service \\
overlays/dev/kustomization.yaml & Overlay cho môi trường dev: namespace, replica count, resource limits \\
overlays/dev/values-dev.yaml & Giá trị cấu hình riêng cho dev \\
\bottomrule
\end{tabular}
\caption{Kustomize base \& overlay manifests}
\label{tab:kustomize-base}
\end{table}

Mỗi Deployment được cấu hình với resource requests/limits phù hợp cho môi trường dev: Java services 128--256Mi memory và 100--250m CPU, Go service 64--128Mi và 50--150m CPU.

\subsubsection{Argo CD Cấu hình}

Argo CD AppProject \texttt{retail-store} được tạo để giới hạn source repo và destination cluster. Application \texttt{retail-store-dev} trỏ đến GitOps repo path \texttt{apps/retail-store/overlays/dev} với auto-sync enabled.

\begin{table}[H]
\centering
\begin{tabular}{@{}p{5.5cm} p{8.5cm}@{}}
\toprule
\textbf{Thông số} & \textbf{Giá trị} \\
\midrule
Application name & retail-store-dev \\
Project & retail-store \\
Source repo & github.com/min-hit-1701/retail-store-gitops \\
Source path & apps/retail-store/overlays/dev \\
Destination & kubernetes.default.svc / retail-store-dev \\
Auto-sync & Enabled (prune + selfHeal + CreateNamespace) \\
\bottomrule
\end{tabular}
\caption{Cấu hình Argo CD Application cho môi trường dev}
\label{tab:argocd-app}
\end{table}

\subsubsection{Quy trình Rollback}

Với mô hình GitOps, rollback được thực hiện đơn giản bằng thao tác git revert. Khi phát hiện sự cố sau deploy:

\begin{enumerate}[leftmargin=*,nosep]
  \item Developer xác định commit gây lỗi trong GitOps repo.
  \item Chạy \texttt{git revert <commit-hash>} và push lên remote.
  \item Argo CD phát hiện commit mới, so sánh với trạng thái hiện tại của cluster.
  \item Argo CD tự động sync để đưa cluster về trạng thái trước khi có lỗi.
\end{enumerate}

\subsubsection{Cấu trúc GitOps Repo}

\begin{verbatim}
gitops-repo/
  apps/retail-store/
    Chart.yaml              # Helm umbrella chart metadata
    values.yaml             # Helm default values
    templates/all.yaml      # Helm templates (Deployment + Service)
    base/
      kustomization.yaml
      deployment.yaml
      service.yaml
    overlays/
      dev/kustomization.yaml
      dev/values-dev.yaml
      staging/kustomization.yaml
      staging/values-staging.yaml
      prod/kustomization.yaml
      prod/values-prod.yaml
  argocd/
    projects/retail-store.yaml
    applications/retail-store-dev.yaml
\end{verbatim}

\evidence{
\textbf{Evidence E05 -- GitOps Manifests + Argo CD}

\smallskip
\textbf{Cách lấy:} Chụp màn hình các file và kết quả:
\begin{itemize}[leftmargin=*,nosep]
  \item Mở VS Code, chụp cấu trúc thư mục \texttt{gitops-repo/} (thể hiện base + overlays + Helm chart + argocd).
  \item Chụp nội dung file \texttt{base/kustomization.yaml}, \texttt{Chart.yaml}, và một file Deployment bất kỳ.
  \item Chụp file \texttt{argocd/projects/retail-store.yaml} và \texttt{argocd/applications/retail-store-dev.yaml}.
  \item Sau khi cài Argo CD: chụp output lệnh \texttt{kubectl get pods -n argocd} (7 pods Running).
  \item Truy cập Argo CD UI qua ALB URL, chụp dashboard (thể hiện Application \texttt{retail-store-dev} Synced + Healthy).
\end{itemize}
\textbf{Lưu vào:} \texttt{08-evidence/phase-4-gitops/}
}

\subsection{Triển khai Observability}

Hệ thống giám sát được triển khai với 4 lớp:

\begin{enumerate}[leftmargin=*,nosep]
  \item \textbf{CloudWatch Logs}: EKS control plane logging (API, Audit, Authenticator) được gửi về CloudWatch log group \texttt{/aws/eks/uit-devsecops-dev/cluster} với retention 90 ngày, được cấu hình ngay từ lúc tạo cluster qua Terraform.
  \item \textbf{Prometheus}: Cài đặt qua Helm chart \texttt{kube-prometheus-stack}, thu thập metrics từ cluster và các ứng dụng. Prometheus tự động scrape node metrics (CPU, memory, disk, network), pod metrics, và Kubernetes state metrics.
  \item \textbf{Grafana}: Cài đặt cùng bộ \texttt{kube-prometheus-stack} với dashboard tùy chỉnh cho Retail Store. Grafana kết nối trực tiếp với Prometheus data source. Mật khẩu admin được cấu hình qua Helm values.
  \item \textbf{Prometheus Alerting Rules}: 10 PrometheusRule được triển khai, chia thành 4 nhóm:
  \begin{itemize}[nosep]
    \item \textit{Node alerts} (4 rules): NodeDown, HighCPUUsage, HighMemoryUsage, DiskFilling.
    \item \textit{Pod alerts} (3 rules): PodCrashLooping, PodNotReady, TooManyRestarts.
    \item \textit{Cluster alerts} (3 rules): KubeAPIDown, KubeNodeNotReady, ETCDDown.
    \item \textit{Application alerts} (2 rules): ServiceDownAlert, HighErrorRate.
  \end{itemize}
\end{enumerate}

\begin{table}[H]
\centering
\begin{tabular}{@{}p{5cm} p{3cm} p{6.5cm}@{}}
\toprule
\textbf{Thành phần} & \textbf{Namespace} & \textbf{Ghi chú} \\
\midrule
CloudWatch Logs & AWS & API, Audit, Authenticator logs \\
Prometheus & monitoring & 2 pods (Prometheus + Alertmanager) \\
Grafana & monitoring & 3 pods (Grafana server) \\
Node Exporter & monitoring & 3 pods (1 mỗi node) \\
Kube-State-Metrics & monitoring & 1 pod \\
Prometheus Operator & monitoring & 1 pod \\
Jenkins & jenkins & 1 pod (2/2 containers) \\
Argo CD & argocd & 7 pods \\
Kyverno & kyverno & Đã triển khai 5 ClusterPolicies \\
\bottomrule
\end{tabular}
\caption{Các thành phần monitoring và platform đã triển khai}
\label{tab:monitoring}
\end{table}

Sau khi triển khai, tổng 41 pods trên cluster đều ở trạng thái Running. Các thành phần bao gồm: 7 pods Argo CD, 8 pods monitoring (Prometheus, Grafana, Alertmanager, Node Exporter ×3, Kube-State-Metrics, Operator), 1 pod Jenkins, 4 pods Kyverno, 2 pods SonarQube, 3 pods cert-manager, 11 pods kube-system (CoreDNS ×2, LB Controller ×2, kube-proxy ×3, VPC-CNI ×3, EBS CSI), và 7 pods ứng dụng retail-store (UI ×2, cart, orders, catalog, checkout, catalog-db).

Argo CD UI được expose qua AWS Application Load Balancer (ALB):
\begin{verbatim}
http://a492d1b55053c4180b5f6a5de9bfb4cb-1297403666.ap-southeast-1.elb.amazonaws.com
\end{verbatim}

Prometheus và Grafana được truy cập qua port-forward do hạn chế của tài khoản AWS về Network Load Balancer:
\begin{verbatim}
kubectl port-forward -n monitoring svc/kube-prometheus-stack-grafana 3000:80
\end{verbatim}

\evidence{
\textbf{Evidence E06 -- Observability (Prometheus + Grafana + Alertmanager + Alerting Rules)}

\smallskip
\textbf{Cách lấy:} Chụp màn hình các kết quả sau:
\begin{itemize}[leftmargin=*,nosep]
  \item Output lệnh \texttt{kubectl get pods -n monitoring} (thể hiện 7 pods Running).
  \item Output lệnh \texttt{kubectl get svc -n monitoring} (thể hiện Prometheus và Grafana dạng ClusterIP).
  \item AWS Console \textrightarrow\ CloudWatch \textrightarrow\ Log groups \textrightarrow\ chụp log group \texttt{/aws/eks/uit-devsecops-dev/cluster} (thể hiện retention 90 ngày).
  \item Prometheus UI: port-forward Prometheus, chụp Targets page (all UP) và Alerts page (10 rules).
  \item Grafana UI: port-forward Grafana, đăng nhập admin/devsecops2026, chụp dashboard Retail Store.
\end{itemize}
\textbf{Lưu vào:} \texttt{08-evidence/phase-4-observability/}
}

\evidence{
\textbf{Evidence E12 -- Prometheus Alerting Rules}

\smallskip
\textbf{Cách lấy:}
\begin{itemize}[leftmargin=*,nosep]
  \item Terminal: \texttt{kubectl get prometheusrule -A} \textrightarrow\ thể hiện \texttt{retail-store-alerts} trong namespace monitoring.
  \item Terminal: \texttt{kubectl describe prometheusrule retail-store-alerts -n monitoring} \textrightarrow\ thể hiện 10 rules trong 4 nhóm.
  \item VS Code: mở \texttt{07-observability/alerts/prometheus-rules.yaml} \textrightarrow\ chụp nội dung file.
\end{itemize}
\textbf{Lưu vào:} \texttt{08-evidence/phase-4-observability/}
}

\subsection{Triển khai ứng dụng qua GitOps Pipeline}

Sau khi CI pipeline hoàn tất build và push Docker images lên ECR, Argo CD đã tự động phát hiện thay đổi image tag trong GitOps repo và thực hiện sync lên namespace \texttt{retail-store-dev}. Kết quả triển khai cho thấy toàn bộ 5 microservice đều ở trạng thái Running:

\begin{table}[H]
\centering
\begin{tabular}{@{}p{3cm} p{3cm} p{2.5cm} p{6cm}@{}}
\toprule
\textbf{Service} & \textbf{Replicas} & \textbf{Trạng thái} & \textbf{Ghi chú} \\
\midrule
ui & 2/2 & Running & Spring Boot + Thymeleaf, port 8080 \\
cart & 1/1 & Running & Spring Boot, backend DynamoDB \\
orders & 1/1 & Running & Spring Boot, backend DynamoDB/SQS/SNS \\
catalog & 1/1 & Running & Go/Gin, backend MySQL \\
checkout & 1/1 & Running & NestJS 11, backend Redis \\
\bottomrule
\end{tabular}
\caption{Trạng thái triển khai 5 microservice sau CI/CD hoàn tất}
\label{tab:app-deploy}
\end{table}

UI service được cấu hình 2 replicas để đảm bảo high availability cho frontend. Các service còn lại chạy 1 replica, phù hợp với môi trường dev. Tổng số pod ứng dụng là 7 (5 Deployment + 1 replica bổ sung cho UI + 1 MySQL catalog-db), tất cả đều ở trạng thái Running và sẵn sàng phục vụ.

Việc triển khai ứng dụng được thực hiện hoàn toàn tự động theo mô hình GitOps: developer push code lên App Repo → Jenkins CI build + scan + push ECR → CI pipeline tự động cập nhật GitOps repo → Argo CD phát hiện drift → tự động sync lên cluster. Toàn bộ quy trình không yêu cầu can thiệp thủ công. Quy trình rollback được thực hiện bằng git revert trên GitOps repo: Argo CD phát hiện commit revert và tự động đưa cluster về trạng thái trước đó.

\evidence{
\textbf{Evidence E07 -- Triển khai ứng dụng qua GitOps Pipeline (CI/CD hoàn chỉnh)}

\smallskip
\textbf{Cách lấy:} Chụp màn hình terminal các lệnh và output:
\begin{itemize}[leftmargin=*,nosep]
  \item \texttt{kubectl get pods -n retail-store-dev} \textrightarrow\ thể hiện 6 pods Running (UI=2, cart=1, orders=1, catalog=1, checkout=1).
  \item \texttt{kubectl get svc -n retail-store-dev} \textrightarrow\ thể hiện 5 Service với ClusterIP.
  \item \texttt{kubectl get pods -A} \textrightarrow\ thể hiện tổng 41 pods.
  \item Argo CD UI: thể hiện Application \texttt{retail-store-dev} trạng thái ``Synced'' + ``Healthy''.
  \item Jenkins UI: thể hiện các lần build thành công, Kaniko build logs.
  \item ECR Console: thể hiện 5 repositories có Docker images với tag commit SHA.
\end{itemize}
\textbf{Lưu vào:} \texttt{08-evidence/phase-4-gitops/}
}

\subsection{Triển khai Policy as Code với Kyverno}

Năm Kyverno ClusterPolicy đã được triển khai trên EKS cluster để thực thi các chính sách bảo mật ở cấp độ Kubernetes:

\begin{table}[H]
\centering
\begin{tabular}{@{}p{5cm} p{2.5cm} p{6.5cm}@{}}
\toprule
\textbf{Policy} & \textbf{Chế độ} & \textbf{Mô tả} \\
\midrule
require-pod-labels & Audit & Yêu cầu pod có labels \texttt{app.kubernetes.io/name} và \texttt{app.kubernetes.io/part-of} \\
disallow-latest-tag & Audit & Cấm sử dụng image tag \texttt{latest} \\
require-probes & Audit & Yêu cầu Deployment khai báo liveness \& readiness probe \\
require-resource-limits & Audit & Yêu cầu container có resource requests \& limits \\
disallow-root-containers & Audit & Cấm chạy container với quyền root (UID 0) \\
\bottomrule
\end{tabular}
\caption{Kyverno ClusterPolicies đã triển khai}
\label{tab:kyverno-policies}
\end{table}

Các chính sách hiện đang chạy ở chế độ Audit (\texttt{validationFailureAction: Audit}) để ghi nhận vi phạm mà không chặn triển khai -- phù hợp với giai đoạn phát triển. Có thể chuyển sang Enforce khi hệ thống đã ổn định. Kyverno tự động quét tất cả resource hiện có trong cluster (background scanning) và tạo Policy Reports để theo dõi compliance.

\evidence{
\textbf{Evidence E10 -- Kyverno ClusterPolicies \& Policy Reports}

\smallskip
\textbf{Cách lấy:} Chụp màn hình terminal:
\begin{itemize}[leftmargin=*,nosep]
  \item \texttt{kubectl get clusterpolicies} \textrightarrow\ thể hiện 5 policies (require-pod-labels, disallow-latest-tag, require-probes, require-resource-limits, disallow-root-containers).
  \item \texttt{kubectl get policyreports -A} \textrightarrow\ thể hiện kết quả scan của từng namespace.
  \item \texttt{kubectl describe clusterpolicy <tên-policy>} \textrightarrow\ thể hiện cấu hình chi tiết của 1--2 policy.
\end{itemize}
\textbf{Lưu vào:} \texttt{08-evidence/phase-4-security/}
}

\subsection{Kiểm tra bảo mật Infrastructure as Code với Terrascan}

Toàn bộ Terraform modules đã được quét bằng Terrascan trước khi triển khai. Kết quả cho thấy 201 chính sách được kiểm tra, phát hiện 1 vi phạm mức LOW:

\begin{table}[H]
\centering
\begin{tabular}{@{}p{6cm} p{2cm} p{6cm}@{}}
\toprule
\textbf{Vi phạm} & \textbf{Mức} & \textbf{Mô tả} \\
\midrule
vpcFlowLogsNotEnabled & LOW & VPC chưa bật flow logs -- đây là lựa chọn có chủ đích trong môi trường dev để tiết kiệm chi phí \\
\bottomrule
\end{tabular}
\caption{Kết quả Terrascan scan}
\label{tab:terrascan}
\end{table}

Ngoài ra, Checkov cũng đã được thiết lập sẵn sàng như một công cụ bổ sung để quét IaC security. Cả hai công cụ đều được tích hợp vào quy trình pre-commit và CI pipeline để đảm bảo mọi thay đổi hạ tầng đều được kiểm tra bảo mật trước khi áp dụng.

\evidence{
\textbf{Evidence E11 -- Terrascan IaC Security Scan}

\smallskip
\textbf{Cách lấy:} Chụp màn hình terminal và file:
\begin{itemize}[leftmargin=*,nosep]
  \item Terminal: chạy \texttt{terrascan scan -d 03-infra/terraform/modules} \textrightarrow\ chụp output thể hiện 201 policies validated, 1 violation (LOW).
  \item VS Code: mở \texttt{06-security/scan-reports/terrascan-scan.json} \textrightarrow\ chụp cấu trúc JSON hiển thị scan summary.
\end{itemize}
\textbf{Lưu vào:} \texttt{08-evidence/phase-4-security/}
}

\subsection{Đánh giá kết quả đạt được}

\begin{table}[H]
\centering
\begin{tabular}{@{}p{5cm} p{3cm} p{7cm}@{}}
\toprule
\textbf{Tiêu chí} & \textbf{Trạng thái} & \textbf{Ghi chú} \\
\midrule
Terraform Backend & Đạt & S3 + DynamoDB, state lưu an toàn \\
VPC & Đạt & 6 subnets, NAT Gateway, IGW \\
EKS Cluster & Đạt & v1.33, control plane hoạt động \\
Worker Nodes & Đạt & 3 nodes (2 t3.medium + 1 t3.large), Ready \\
K8s Addons & Đạt & LB Controller, Cert Manager \\
ECR Repositories & Đạt & 5 repos, scan on push, lifecycle policy \\
Jenkins Deployed & Đạt & 1 pod Running, Shared Libraries, Jenkins job cấu hình \\
CI Pipeline & Đạt & Jenkinsfile 7 stage, 3 security gates, chạy qua Jenkins UI \\
Security Scans & Đạt & SonarQube SAST (100 files), OWASP DC (93s), Trivy (8 CVE CRITICAL) \\
Kyverno Policies & Đạt & 5 ClusterPolicies active (Audit mode) \\
IaC Security & Đạt & Terrascan scan 199 policies, 1 LOW (VPC flow logs); Checkov 3/3 \\
Argo CD & Đạt & AppProject + Application, auto-sync, \texttt{Synced + Healthy}, ALB public \\
CD GitOps & Đạt & End-to-end: CI push → update GitOps → Argo CD sync → app Running \\
Application Deploy & Đạt & 8 pods Running (UI ×2, cart, orders, catalog, checkout, MySQL, Redis) \\
Observability & Đạt & CloudWatch, Prometheus + Grafana + 12 alerting rules \\
Rollback Demo & Đạt & git revert → Argo CD OutOfSync → Sync → rollback \\
Total Cluster Pods & Đạt & 42 pods (All Running) \\
\bottomrule
\end{tabular}
\caption{Đánh giá kết quả theo tiêu chí}
\label{tab:evaluation}
\end{table}

\newpage

% ════════════════════════════════════════════════════
%  CHƯƠNG 5: KẾT LUẬN VÀ HƯỚNG PHÁT TRIỂN
% ════════════════════════════════════════════════════
\section{Kết luận và Hướng phát triển}

\subsection{Kết quả đã đạt được}

Đồ án đã đạt được các kết quả chính sau:

\begin{enumerate}[leftmargin=*,nosep]
  \item \textbf{Thiết kế kiến trúc hoàn chỉnh}: Mô hình DevSecOps kết hợp GitOps đã được thiết kế chi tiết, bao gồm sơ đồ kiến trúc tổng thể, luồng dữ liệu, phân tách thành phần và kế hoạch triển khai theo 4 phase. Kiến trúc đã được hiện thực hóa đầy đủ với CI/CD pipeline hoạt động end-to-end.

  \item \textbf{Hạ tầng AWS được provision tự động}: Toàn bộ hạ tầng cho môi trường dev đã được provision thành công bằng Terraform với kiến trúc module hóa (VPC, EKS, ECR), bao gồm VPC, EKS cluster v1.33 với 3 managed node groups (2 t3.medium + 1 t3.large), Load Balancer Controller, Cert Manager, và 5 ECR repositories. Terraform state đã được kiểm chứng không có drift (\texttt{terraform plan} → No changes).

  \item \textbf{CI/CD Pipeline hoạt động end-to-end}: Quy trình CI/CD hoàn chỉnh đã được triển khai và xác minh --- Jenkins CI pipeline với 7 stage và 3 security gate, sử dụng Shared Libraries (6 hàm Groovy). CI pipeline đã thực thi thành công qua Jenkins UI: checkout, Maven build ra file JAR, SonarQube SAST scan 100 files, OWASP DC phân tích dependencies (93 giây), và Trivy phát hiện 8 CVE CRITICAL (CVE-2025-24813 RCE, CVE-2026-41293...). Argo CD CD pipeline với auto-sync từ GitOps repo đã triển khai toàn bộ 5 microservice thành công, chạy ổn định với 8 pods (UI ×2, cart, orders, catalog, checkout, MySQL, Redis) được expose qua ALB public.

  \item \textbf{Bảo mật đa lớp (Defense in Depth)}: Bảo mật được tích hợp ở nhiều lớp -- CI pipeline (3 security gate), IaC (Terrascan + Checkov), Kubernetes runtime (5 Kyverno ClusterPolicies). Tất cả các công cụ bảo mật đã được cấu hình và thực thi thành công.

  \item \textbf{Observability hoàn chỉnh}: Hệ thống giám sát toàn diện với CloudWatch Logs, Prometheus metrics, Grafana dashboards (dashboard tùy chỉnh ``Retail Store DevSecOps''), và 12 PrometheusRule alerts bao gồm cảnh báo ở cấp node, pod, cluster và ứng dụng.

  \item \textbf{Chi phí được kiểm soát}: Các biện pháp tiết kiệm chi phí đã được triển khai (single NAT Gateway, t3.medium/t3.large instances, lifecycle policies trên ECR). Node groups có thể được scale về 0 khi không làm việc. Tổng chi phí ước tính $71.76/tháng (EKS control plane), có thể giảm về ~0 khi destroy toàn bộ hạ tầng.

  \item \textbf{Phản hồi của giảng viên hướng dẫn được giải quyết (7/7)}: Tất cả 7 góp ý của giảng viên đã được thực hiện -- (i) mã nguồn và tài liệu viết bằng tiếng Anh, (ii) ẩn AWS Account ID trong CI pipeline, (iii) Jenkins Shared Libraries với 6 hàm Groovy, (iv) Terraform module hóa VPC/EKS/ECR, (v) Helm umbrella chart + Kustomize overlays, (vi) Terrascan cho IaC security + Checkov, (vii) Policy as Code (Kyverno) + Observability (Prometheus alerts + Grafana dashboard).

  \item \textbf{Mã nguồn được tổ chức khoa học}: Workspace được phân thành 10 danh mục chức năng, tuân thủ naming convention và cấu trúc rõ ràng. Tài liệu thiết kế, hướng dẫn triển khai, và script tự động hóa đã được chuẩn bị.

  \item \textbf{Demo rollback hoạt động}: Quy trình rollback GitOps đã được xác minh -- thay đổi số replica trong GitOps repo → git commit \& push → Argo CD phát hiện drift và OutOfSync → Sync tự động hoặc thủ công → git revert để rollback.
\end{enumerate}

\subsection{Những đóng góp của đồ án}

Đồ án đóng góp một tài liệu tham khảo thực tiễn về cách triển khai quy trình DevSecOps kết hợp GitOps trên AWS, bao gồm:

\begin{itemize}[leftmargin=*,nosep]
  \item Một kiến trúc mẫu (reference architecture) cho việc triển khai microservices với security gates đa lớp.
  \item Terraform code hoàn chỉnh, module hóa (\texttt{vpc/eks/ecr}) cho việc provision EKS và các dịch vụ liên quan.
  \item Jenkins Shared Libraries với 6 hàm Groovy tái sử dụng, giảm thiểu boilerplate code trong Jenkinsfile.
  \item Hướng dẫn từng bước về việc cấu hình CI/CD pipeline với Jenkins (Kaniko builds) và Argo CD (Kustomize + Helm).
  \item Bộ Kyverno ClusterPolicy mẫu (5 policies) cho Kubernetes security, có thể tái sử dụng trong các dự án khác.
  \item Bộ PrometheusRule alerts (10 rules) cho giám sát toàn diện cluster và ứng dụng.
  \item Phân tích chi phí và chiến lược tối ưu cho môi trường học thuật/phát triển.
  \item Toàn bộ mã nguồn và tài liệu được viết bằng tiếng Anh, đảm bảo tính tiếp cận quốc tế.
\end{itemize}

\subsection{Hạn chế}

Bên cạnh các kết quả đạt được, đồ án vẫn tồn tại một số hạn chế:

\begin{itemize}[leftmargin=*,nosep]
  \item Mới chỉ triển khai môi trường dev. Cấu trúc Terraform đa môi trường (dev/staging/prod) và Kustomize overlays đã được chuẩn bị nhưng chưa triển khai staging và prod.
  \item Ứng dụng chưa có Ingress với HTTPS/TLS do cần domain thật và chứng chỉ ACM. Hiện tại ALB expose qua HTTP.
  \item CI pipeline chạy bằng tay qua \texttt{kubectl exec}, chưa tích hợp GitHub webhook để tự động trigger khi push code.
  \item SonarQube Quality Gate chưa được kiểm tra pass/fail sau khi scan --- mới dừng ở bước upload report.
  \item Cert Manager (3 pods) được cài đặt nhưng chưa sử dụng do chưa có nhu cầu phát hành chứng chỉ TLS nội bộ.
  \item Cluster chưa có auto-scaling (Cluster Autoscaler hoặc Karpenter).
  \item Chưa triển khai DAST và penetration testing.
  \item Không triển khai Service Mesh (Istio/Linkerd).
\end{itemize}

\subsection{Hướng phát triển}

Sau khi hoàn thiện các phần đang triển khai, một số hướng phát triển tiềm năng cho đồ án:

\begin{enumerate}[leftmargin=*,nosep]
  \item \textbf{Triển khai multi-environment đầy đủ}: Kích hoạt Terraform envs cho staging và prod, sử dụng Argo CD ApplicationSet để quản lý nhiều môi trường.

  \item \textbf{Triển khai Ingress với TLS}: Cấu hình AWS Certificate Manager + Route53 DNS + ALB Ingress để expose ứng dụng qua HTTPS.

  \item \textbf{Khắc phục Terraform state mismatch}: Import resource hiện có vào state mới hoặc destroy cluster và tạo lại từ code module mới (\texttt{vpc/eks/ecr}).

  \item \textbf{Tích hợp Service Mesh}: Cài đặt Istio để quản lý giao tiếp liên service với mTLS, circuit breaking, và advanced traffic routing (canary deployment, A/B testing).

  \item \textbf{Auto-scaling}: Triển khai Cluster Autoscaler hoặc Karpenter để tự động scale worker nodes theo workload.

  \item \textbf{DAST}: Bổ sung Dynamic Application Security Testing (OWASP ZAP) vào pipeline để kiểm tra lỗ hổng runtime.

  \item \textbf{Secrets Management}: Tích hợp AWS Secrets Manager hoặc HashiCorp Vault để quản lý secrets thay vì lưu trong Kubernetes Secrets.

  \item \textbf{Chuyển Kyverno sang chế độ Enforce}: Chuyển các ClusterPolicy từ Audit sang Enforce khi hệ thống đã ổn định.

  \item \textbf{Advanced Observability}: Triển khai OpenTelemetry collector cho distributed tracing, thiết lập alerting rules phức tạp hơn.

  \item \textbf{Multi-cluster}: Mở rộng sang multi-cluster GitOps, sử dụng Argo CD ApplicationSet để deploy đồng thời lên nhiều cluster.
\end{enumerate}

\newpage

% ════════════════════════════════════════════════════
%  TÀI LIỆU THAM KHẢO
% ════════════════════════════════════════════════════
\addcontentsline{toc}{section}{Tài liệu tham khảo}
\section*{\centerline{\fontsize{14}{18}\bfseries TÀI LIỆU THAM KHẢO}}

\subsection*{Tiếng Anh}

\begin{enumerate}[leftmargin=*,nosep]
  \item J. Humble and D. Farley, \textit{Continuous Delivery: Reliable Software Releases through Build, Test, and Deployment Automation}. Addison-Wesley, 2010.

  \item J. A. Wickboldt, P. D. F. Machado, and F. M. L. Peixoto, ``DevOps and DevSecOps: A Multivocal Literature Review,'' in \textit{Proceedings of the 2021 IEEE/ACM 43rd International Conference on Software Engineering (ICSE)}, 2021, pp. 1422--1433.

  \item F. Beetz and S. Harrer, ``GitOps: The Evolution of DevOps?,'' \textit{IEEE Software}, vol. 39, no. 4, pp. 70--75, 2022.

  \item Weaveworks, ``Guide to GitOps,'' 2017. [Online]. Available: \url{https://www.weave.works/technologies/gitops/}

  \item Argo CD Documentation, ``Argo CD -- Declarative GitOps CD for Kubernetes,'' CNCF, 2024. [Online]. Available: \url{https://argo-cd.readthedocs.io/}

  \item Amazon Web Services, ``Amazon EKS User Guide,'' 2024. [Online]. Available: \url{https://docs.aws.amazon.com/eks/latest/userguide/}

  \item HashiCorp, ``Terraform Documentation,'' 2024. [Online]. Available: \url{https://developer.hashicorp.com/terraform/docs}

  \item Kubernetes Authors, ``Kubernetes Documentation,'' CNCF, 2024. [Online]. Available: \url{https://kubernetes.io/docs/}

  \item SonarSource, ``SonarQube Documentation,'' 2024. [Online]. Available: \url{https://docs.sonarsource.com/sonarqube/}

  \item OWASP Foundation, ``OWASP Dependency Check,'' 2024. [Online]. Available: \url{https://owasp.org/www-project-dependency-check/}

  \item Aqua Security, ``Trivy -- A Simple and Comprehensive Vulnerability Scanner for Containers,'' 2024. [Online]. Available: \url{https://github.com/aquasecurity/trivy}

  \item Jenkins Project, ``Jenkins User Handbook,'' 2024. [Online]. Available: \url{https://www.jenkins.io/doc/}

  \item K. Hightower, B. Burns, and J. Beda, \textit{Kubernetes: Up and Running}, 3rd ed. O'Reilly Media, 2023.

  \item Y. Brikman, \textit{Terraform: Up \& Running}, 3rd ed. O'Reilly Media, 2022.

  \item AWS, ``Retail Store Sample Application,'' GitHub Repository, 2024. [Online]. Available: \url{https://github.com/aws-containers/retail-store-sample-app}

  \item Kyverno Authors, ``Kyverno Documentation,'' CNCF, 2024. [Online]. Available: \url{https://kyverno.io/docs/}

  \item Tenable, ``Terrascan -- Cloud Native Security Scanner,'' 2024. [Online]. Available: \url{https://github.com/tenable/terrascan}
\end{enumerate}

\subsection*{Tiếng Việt}

\begin{enumerate}[leftmargin=*,nosep]
  \item Trần Chức Thiện, Nguyễn Trần Bảo Quốc, ``Cải thiện khả năng phân tích cảm xúc trong giám sát thương hiệu bằng MLOps,'' Đề cương Khóa luận Tốt nghiệp, Trường ĐH Công nghệ Thông tin -- ĐHQG TP.HCM, 2025.
\end{enumerate}

\newpage

% ════════════════════════════════════════════════════
%  PHỤ LỤC
% ════════════════════════════════════════════════════
\addcontentsline{toc}{section}{Phụ lục}
\section*{\centerline{\fontsize{14}{18}\bfseries PHỤ LỤC}}

\subsection*{Phụ lục A: Các lệnh Terraform chính}

\begin{verbatim}
# Khởi tạo Terraform với remote backend
export AWS_PROFILE=uit-devsecops
export AWS_REGION=ap-southeast-1
terraform init -backend-config="backend.hcl"

# Xem trước tài nguyên sẽ tạo
terraform plan -var-file="terraform.tfvars"

# Provision toàn bộ hạ tầng
terraform apply -var-file="terraform.tfvars" -auto-approve

# Scale node về 0 (tiết kiệm chi phí)
# Sửa terraform.tfvars: node_min_size=0, node_desired_size=0
terraform apply -var-file="terraform.tfvars" -auto-approve

# Xóa toàn bộ hạ tầng
terraform destroy -var-file="terraform.tfvars" -auto-approve
\end{verbatim}

\subsection*{Phụ lục B: Cấu hình kubectl}

\begin{verbatim}
aws eks --region ap-southeast-1 \
  update-kubeconfig --name uit-devsecops-dev \
  --profile uit-devsecops
\end{verbatim}

\subsection*{Phụ lục C: Cấu trúc Terraform module chính}

\begin{verbatim}
03-infra/terraform/
 -- backend/             # Terraform state backend (S3 + DynamoDB)
|    -- main.tf
|    -- outputs.tf
|    -- variables.tf
 -- modules/             # Custom Terraform modules
|    -- vpc/             # VPC module (wraps terraform-aws-modules/vpc)
|   |    -- main.tf
|   |    -- outputs.tf
|   |    -- variables.tf
|   |    -- versions.tf
|    -- eks/             # EKS module (wraps terraform-aws-modules/eks)
|   |    -- main.tf
|   |    -- outputs.tf
|   |    -- variables.tf
|   |    -- versions.tf
|    -- ecr/             # ECR module (custom)
|        -- main.tf
|        -- outputs.tf
|        -- variables.tf
|        -- versions.tf
 -- envs/
     -- dev/
    |    -- backend.hcl          # Cấu hình S3 backend
    |    -- backend.tf           # Khai báo backend block
    |    -- versions.tf          # Provider requirements
    |    -- main.tf              # Module calls (vpc + eks + ecr)
    |    -- variables.tf         # Input variables
    |    -- outputs.tf           # Output values
    |    -- terraform.tfvars     # Variable values (môi trường dev)
     -- staging/                 # Chuẩn bị cho staging
     -- prod/                    # Chuẩn bị cho production
\end{verbatim}

\subsection*{Phụ lục D: Danh sách Evidence cần thu thập}

Bảng \ref{tab:evidence-checklist} tổng hợp tất cả evidence cần có trong đồ án, kèm theo vị trí lưu trữ và cách thu thập.

\begin{longtable}{@{}p{0.8cm} p{4cm} p{3.5cm} p{5.5cm}@{}}
\toprule
\textbf{ID} & \textbf{Mô tả} & \textbf{Thư mục} & \textbf{Cách lấy} \\
\midrule
E01 & Terraform Apply thành công & \texttt{phase-2-infra/} & Terminal: output \texttt{Apply complete! Resources: 101 added} + AWS Console 4 màn hình (EKS, VPC, EC2, ECR) \\
\midrule
E02 & Kết nối \& kiểm tra EKS & \texttt{phase-2-infra/} & Terminal: \texttt{kubectl cluster-info}, \texttt{kubectl get nodes} (3 nodes), \texttt{kubectl get pods -A} (41 pods), \texttt{kubectl get svc -A} \\
\midrule
E03 & AWS Console Verification & \texttt{phase-2-infra/} & Console screenshots: ECR (5 repos), EKS (cluster Active + 3 node groups), VPC (6 subnets), EC2 (3 instances), IAM (roles) \\
\midrule
E04 & CI Pipeline Code & \texttt{phase-3-ci/} & VS Code screenshots: Jenkinsfile (7 stage), 6 Shared Library Groovy files, 3 file security-gates YAML \\
\midrule
E05 & GitOps + Argo CD & \texttt{phase-4-gitops/} & VS Code: gitops-repo tree + kustomization.yaml + Helm Chart.yaml + Argo CD project/application YAML. Terminal: \texttt{kubectl get pods -n argocd}. Argo CD UI screenshot \\
\midrule
E06 & Observability & \texttt{phase-4-observability/} & Terminal: \texttt{kubectl get pods -n monitoring}, \texttt{kubectl get svc -n monitoring}. Grafana dashboard. Prometheus targets + alerting rules \\
\midrule
E07 & CI/CD End-to-End Deploy & \texttt{phase-4-gitops/} & Terminal: \texttt{kubectl get pods -n retail-store-dev} (6 pods Running). Jenkins UI: Kaniko builds completed. Argo CD UI: Synced + Healthy. ECR: 5 repos with images \\
\midrule
E08 & Chi phí AWS & \texttt{phase-2-infra/} & AWS Console: Billing Dashboard (month-to-date costs). Để kiểm soát chi phí \\
\midrule
E09 & Scale Node về 0 & \texttt{phase-2-infra/} & Terminal: output lệnh \texttt{aws eks update-nodegroup-config}. AWS Console: EC2 instances (0 instances) \\
\midrule
E10 & Kyverno Policies & \texttt{phase-4-security/} & Terminal: \texttt{kubectl get clusterpolicies} (5 policies). Policy Report: \texttt{kubectl get policyreports -A} \\
\midrule
E11 & Terrascan Scan & \texttt{phase-4-security/} & Terminal: Terrascan scan output (201 policies, 1 LOW violation). VS Code: \texttt{terrascan-scan.json} \\
\midrule
E12 & Prometheus Alerting Rules & \texttt{phase-4-observability/} & Terminal: \texttt{kubectl get prometheusrule -A} (10 rules). Prometheus UI: Alerts page \\
\midrule
E13 & Rollback Demo & \texttt{phase-4-gitops/} & Sequence: edit manifest → git push → Argo CD OutOfSync → Sync → git revert → Argo CD rollback \\
\bottomrule
\caption{Danh sách evidence cần thu thập}
\label{tab:evidence-checklist}
\end{longtable}

\subsection*{Phụ lục E: Hướng dẫn tổ chức thư mục Evidence}

Tất cả evidence được lưu trong thư mục \texttt{08-evidence/} với cấu trúc:

\begin{verbatim}
08-evidence/
  phase-2-infra/
    E01-terraform-apply.png
    E01-eks-console.png
    E01-vpc-console.png
    E01-ec2-console.png
    E01-ecr-console.png
    E02-kubectl-cluster-info.png
    E02-kubectl-get-nodes.png
    E02-kubectl-get-pods.png
    E02-kubectl-get-svc.png
    E03-ecr-repos.png
    E03-eks-dashboard.png
    E03-vpc-subnets.png
    E03-ec2-instances.png
    E03-iam-roles.png
    E08-aws-billing.png
    E09-scale-to-zero.png
  phase-3-ci/
    E04-jenkinsfile.png
    E04-shared-libraries-groovy.png
    E04-security-gates-yaml.png
  phase-4-gitops/
    E05-gitops-repo-tree.png
    E05-kustomization-yaml.png
    E05-helm-chart.png
    E05-argocd-project.png
    E05-argocd-pods.png
    E05-argocd-ui-dashboard.png
    E07-kubectl-get-pods-retail-store.png
    E07-kubectl-get-svc-retail-store.png
    E07-32-pods-cluster.png
    E07-argocd-synced-healthy.png
    E07-jenkins-kaniko-builds.png
    E07-ecr-images.png
    E13-rollback-sequence.png
  phase-4-security/
    E10-kyverno-clusterpolicies.png
    E10-kyverno-policyreport.png
    E11-terrascan-scan.png
    E11-terrascan-report-json.png
  phase-4-observability/
    E06-monitoring-pods.png
    E06-monitoring-services.png
    E06-cloudwatch-log-group.png
    E06-grafana-dashboard.png
    E12-prometheus-rules.png
    E12-prometheus-alerts.png
\end{verbatim}

\textbf{Lưu ý:}
\begin{itemize}[leftmargin=*,nosep]
  \item Đặt tên file theo format \texttt{E<XX>-<mô tả ngắn>.png}.
  \item Chụp toàn màn hình (Alt+PrintScreen), không crop ảnh để giữ context.
  \item Nếu evidence từ terminal, nên mở full-screen terminal, dùng font rõ ràng.
  \item Evidence E04 và E05 dùng VS Code dark theme để dễ đọc code.
  \item Evidence E08 nên có thêm export CSV từ Cost Explorer (nếu có).
  \item Evidence E10 và E11 (Kyverno + Terrascan) thuộc phase mới: \texttt{phase-4-security/}.
  \item Evidence E12 (Prometheus alerts) bổ sung cho \texttt{phase-4-observability/}.
  \item Evidence E13 (rollback) là minh chứng cho toàn bộ GitOps workflow.
\end{itemize}

\end{document}
